\documentclass[../main.tex]{subfiles}
\begin{document}
%==========================================TIÊU ĐỀ TẬP========================================
\begin{table}[h]
  \begin{adjustwidth}{-1cm}{}
    \begin{tabular}{>{\centering\arraybackslash}p{6cm}|>{\raggedright\arraybackslash}p{12cm}}
      \multirow{5}{6cm}
      % Cột bên trái: ÉPISODE và số tập
      {\\[-40pt]
        \centering
        {\fontsize{20pt}{20pt}\selectfont{\outlinetext[1]{eptype}{white}{\flaregothic PERSONNALISÉ}}\\[12pt]      % use for normal/special episodes
          {\fontsize{36pt}{36pt}\selectfont{\outlinetext[1]{epnum}{white}{\burbank BẢY}}}\\[6pt]     % use for custom episodes
          {\fontsize{20pt}{20pt}\selectfont{\outlinetext[1]{epnum}{white}{\cafeteria (C-7)}}}\color{black}}
        \\[18pt]
        % {\fontsize{14pt}{14pt}\selectfont{\outlinetext[1]{epattr}{white}{\flaregothic AU REVOIR, }}{\flaregothic\color{Rushineko} Rushia}\fontsize{14pt}{14pt}\selectfont{\outlinetext[1]{epattr}{white}{\flaregothic \ !}}}\\
        % {\fontsize{18pt}{18pt}\selectfont{\outlinetext[1]{epattr}{white}{\flaregothic FIN DE TOME !}}}\\
        \color{black}
      }
       &
      % Dòng thứ nhất: tên tập tiếng Nhật
      {\fotskip\fontsize{18pt}{18pt}\selectfont \ruby{平}{へい}\ruby{和}{わ}\ruby{的}{てき}な\ruby{終}{お}わる、るしあ……}                                  \\[6pt]
      % Dòng thứ hai: tên tập tiếng Anh
       & {\cafeteria\fontsize{18pt}{18pt}\selectfont What-If: Rushia's Good Ending?}                                                      \\[4pt]
      % Dòng thứ ba: tên tập tiếng Việt
       & {\vnmsans\fontsize{18pt}{18pt}\selectfont Buổi tốt nghiệp trong mơ}                                                              \\[8pt]
      % Dòng thứ tư: Nhân vật xuất hiện
       & {\caviar{\fontsize{14pt}{14pt}\selectfont Nhân vật: \emp{kuon2} \emp{pekora} \emp{rushia} \emp{flare1} \emp{noel} \emp{marine}}} \\[8pt]
      % Dòng cuối cùng: Thời gian diễn ra của tập (thường sẽ là ngày phát hành)
       & {\sen\fontsize{16pt}{16pt}\selectfont Mốc thời gian: 22/02/2022}                                                                 \\[18pt]
    \end{tabular}\\[8pt]
  \end{adjustwidth}
\end{table}
%=========================================TRUYỆN CHÍNH========================================
\color{black}

\txtbx[2]{
  {\color{vol} \uline{Tóm tắt:}} Kuon mơ thấy mình quay lại thời cấp ba, lớp 12A3 đón năm thành viên Gen 3 (Pekora, Rushia, Flare, Noel, Marine) vào buổi học cuối cùng. Từ tiết Toán, Lý, Hóa đến Tiếng Anh, mỗi người bộc lộ tài năng (Flare giỏi Hình, Pekora lý ``cân team'', Rushia hóa xuất sắc) và những pha hài (Noel khoanh đại, Marine đọc tiếng Anh sai). Buổi sinh hoạt cuối biến thành lễ tốt nghiệp cho Rushia với những lời tri ân xúc động. Tỉnh dậy, Kuon nhận tin Rushia mất tích sau khi bị sa thải, giấc mơ về một ``kết thúc tốt đẹp'' cho nàng Pháp sư.
}

\subsubsection*{\phantomsection\label{ep-C7.1}}
\addcontentsline{toc}{subsubsection}{{\sen{-Partie 1-}} {\caviar Buổi học đầu tiên (và cuối cùng)}}
\stpt{-Partie 1-}\stcap{Buổi học đầu tiên (và cuối cùng)}

\pov{Hikawa Kuon}

\lang{Etsunan}

``\dots đi! Dậy đi xem nào!''

Một giọng nói gọi giật ngược kéo tôi ra khỏi tầng tầng lớp lớp hư vô. Tôi uể oải rên rỉ một tiếng trong cổ họng, định bụng xoay đầu sang hướng khác để chìm tiếp vào giấc ngủ. Nhưng tiếng đập bàn tàn nhẫn vang lên ngay sát tai đã đập tan ý định đó.

``Dậy đi ông cháu ơi. Đến giờ truy bài rồi đấy!'' một giọng nói khác có vẻ nhắng nhít hơn lại vang lên từ tai bên kia.

``\dots Dở người à\dots'' tôi lầm bầm, giọng khàn đặc vì ngái ngủ. ``Lại trêu tôi chứ gì\dots''

``Trêu gì mà trêu. Trống vào lớp đánh từ đời nào rồi bạn.''

``\dots Hả?''

Tôi chậm rãi ngẩng mặt lên. Đôi mắt còn mờ sương nhấp nháy liên tục để điều tiết ánh sáng.

Trong một khoảnh khắc ngắn ngủi, một luồng khí lạnh buốt bất ngờ chạy dọc sống lưng làm tôi giật thót. Trí óc tôi vô thức lóe lên hình ảnh những hành lang quánh đặc sương mù, mùi lá mục nồng nặc gắt gỏng, cùng những lớp học hoang tàn của trường Trung học Aogami, nơi tôi, Kaigou và Suzuki từng phải vật vã, kiệt sức trong những vòng lặp sinh tử đầy vô vọng. Cái cảm giác bị giam cầm trong ngôi trường cấp ba tàn khốc đó chỉ vừa mới trôi qua chưa lâu, khiến lồng ngực tôi bất giác thắt lại, nhịp tim đập dồn dập vì ngột ngạt.

Nhưng rồi, sự tỉnh táo chậm rãi quay trở lại.

Nơi này hoàn toàn không có không khí u uất hay bóng dáng tàn khốc nào của lời nguyền.

Khung cảnh hiện ra trước mắt lại là một nơi vô cùng quen thuộc kể từ khi tôi tốt nghiệp cấp ba. Chiếc bảng xanh bám đầy vết phấn trắng, những chiếc bàn gỗ đôi xếp thành bốn dãy thẳng tắp, cùng bức tường sơn vàng đã bị nứt dăm vài chỗ do vết mòn thời gian. Tuy cổ kính là thế, phòng học vẫn được trang bị vài món công nghệ khá hiện đại: chiếc TV màn hình lớn lắp ngay sau bảng xanh, hai cái điều hòa chạy êm ru cùng bốn chiếc quạt trần vừa thay mới đang quay đều đều.

Cúi đầu nhìn xuống, tôi nhận ra mình đang mặc bộ đồng phục của trường cấp ba cũ - Trường THPT \ruby{Anwa}{Yên Hòa}, một trong những ngôi trường điểm của quận Kyuukami nói riêng và toàn Kawachi nói chung. Đó là chiếc áo phông trắng có cổ viền xanh đậm, với logo nhà trường đính trang trọng trên tấm vải trước ngực trái, kết hợp cùng quần tây xám thanh lịch. Tôi vẫn đi chiếc dép quai hậu quen thuộc.

Nhìn quanh, tôi nhận ra ba người bạn đang vây quanh bàn mình với những biểu cảm vô cùng đa dạng. Zetaro và Khan là hai người đứng sát nhất.

``Dậy rồi sao, anh bạn?''

Một giọng nói điềm nhiên cất lên từ dãy bàn bên cạnh.

Satago Hakimi (\ruby{佐}{さ}\ruby{太}{た}\ruby{語}{ご}\ruby{葉}{は}\ruby{君}{きみ}, \ruby{Xa}{Tá}-\ruby{Thái}{ta}-\ruby{Ngự}{gô}\ \ruby{Ha}{Diệp}-\ruby{ki-mi}{Quân}). Cậu chàng học chung lớp với tôi từ năm lớp mười một. Hakimi có trình độ tiếng Anh giỏi đến mức từng ẵm giải cao trong kỳ thi học sinh giỏi của trường, một thành tích mà một đứa chẳng mảy may quan tâm đến học thuật như tôi chỉ biết đứng nhìn. Cậu ta tự nhận tính mình khá cục súc. Ở thế giới thực, nếu tôi nhớ không nhầm, giờ này Hakimi đang đi du học bên Úc rồi.

Gọi là bạn thân thì cũng không hẳn, vì chúng tôi mới chỉ biết nhau hai năm, và sau khi tốt nghiệp cấp ba, hầu như chúng tôi cũng không trao đổi qua lại nhau là mấy. Nói rằng ``đôi bạn cùng chung hoàn cảnh'' có lẽ sẽ hợp lý hơn, bởi cả tôi và cậu ta đều là học sinh từ lớp khác chuyển sang, tính tình lại thuộc dạng trầm lặng như tôi.

\dots Nhưng khoan đã. Tại sao tôi lại ở đây? Tôi tưởng mình đã tốt nghiệp đại học đến nơi rồi cơ mà?

``\dots Tôi đang ở trong lớp 12A3 đấy à?'' tôi ngơ ngác hỏi.

``Chứ còn sao nữa. Tỉnh ngủ đi bạn,'' Hakimi đáp, gương mặt không một gợn sóng.

``Tôi\dots\ ngủ được bao lâu rồi?''

Khan chêm vào: ``Ông ngủ được tầm nửa tiếng đấy.''

``Vãi\dots''

Hakimi, Khan, Zetaro và tôi chơi chung trong một hội bạn mà cả lũ vẫn hay tự gọi bằng cái tên cực kỳ thô bựa: nhóm ``Rách'' (viết tắt của một cụm từ tục tĩu mà chẳng ai buồn nhắc lại). Cho đến tận bây giờ, tôi vẫn không hiểu tại sao cái nhóm này lại đặt một cái tên rẻ tiền đến thế.

Còn lý do tại sao một kẻ lười giao thiệp như tôi lại có mặt trong đó? Theo như Zetaro từng kể, là bởi vì thấy tôi lúc nào cũng cô độc một mình nên bọn họ mới chủ động kéo tôi vào hội. Đối với một đứa ưa sự yên tĩnh nhưng lại thuộc tuýp ``sao cũng được'' như tôi, tôi chẳng hề phản đối. ``Bọn mày muốn thế nào chả được,'' tôi đã nói thế vào cái ngày gia nhập.

Và rồi bùm một cái, chúng tôi chơi thân với nhau từ lúc nào không hay. Ngoài bốn đứa tôi ra, nhóm ``Rách'' còn có thêm vài gương mặt cộm cán khác:

Koheki, cậu bạn giỏi Hóa tới mức thần đồng mà tôi từng kể sau cái vụ ``quả bom\footnote{{\flaregothic ÉPISODE 62}, quyển số Năm của \textit{holo no Graffiti}.}'' hồi đó. Cậu ta cực kỳ thích ngồi bàn cuối. Cứ đến giờ Hóa là Koheki lại lôi việc riêng ra làm, nhưng thầy giáo bộ môn cũng đành tặc lưỡi cho qua vì đơn giản là cậu ta giỏi đến mức không cần nghe giảng nữa.

Todorote Boji (\ruby{轟}{とどろ}\ruby{天}{テン}\ruby{慕}{ぼ}\ruby{情}{じ}, \ruby{Tô-dô-rô}{Oanh}-\ruby{tê}{Thiên}\ \ruby{Bô}{Mộ}-\ruby{gi}{Tình}) - một gã tính tình vô cùng dễ dãi. Sức học trên lớp của cậu ta vốn chỉ ngang ngửa tôi, nhưng không hiểu sao đến kỳ thi đại học lại đạt số điểm cao chót vót. Tất nhiên tôi không ghen tị, nhưng đó vẫn là một ẩn số mà sau từng ấy năm tôi vẫn chưa tài nào giải mã nổi. Boji còn có biệt danh là ``Invisible'' (Kẻ vô hình) vì gần như chẳng bao giờ ngoi lên nhắn tin trong nhóm chat trực tuyến.

Rutobaga Kijiru (\ruby{劉}{る}\ruby{登}{と}\ruby{簿}{ば}\ruby{外}{が}\ruby{気}{き}\ruby{汁}{じる}, \ruby{Rư}{Lưu}-\ruby{tô}{Đăng}-\ruby{ba}{Bạc}-\ruby{ga}{Ngoại}\ \ruby{Ki}{Khí}-\ruby{gi-rư}{Hiệp}) - cậu bạn có vóc dáng khá khiêm tốn nhưng bù lại cực kỳ giỏi các môn khoa học tự nhiên. Dù sao thì lớp chúng tôi cũng là lớp chuyên Toán - Lý mà. Kijiru rất xông xênh, hay bao tôi ăn uống mỗi lần tôi lỡ quên mang tiền.

Motoyoga Kitaharu (\ruby{元}{も}\ruby{々}{も}\ruby{洋}{よう}\ruby{楽}{が}\ruby{北}{きた}\ruby{治}{はる}, \ruby{Mô}{Nguyên}-\ruby{tô}{Nguyên}-\ruby{giô}{Dương}-\ruby{ga}{Lạc}\ \ruby{Ki-ta}{Bắc}-\ruby{ha-rư}{Trị}) - cạ cứng của Zetaro. Cậu ta cuồng các trò chơi điện tử trực tuyến, chủ yếu là ``Li*n Qu*n'' và thi thoảng là ``Lửa Chùa''. Nghe đồn Kitaharu rất thích cắm rốt ngoài quán net để leo rank cùng cái gã Zetaro dâm đặng ấy. Theo tôi nhớ, về sau hai người đó học cùng trường Bách khoa, nằm ngay sát vách trường Xây dựng của tôi.

Tóm lại, trừ tôi và Boji thuộc dạng trung bình tầm trung ra, những thành viên còn lại trong nhóm đều là học sinh có số có má về mảng học tập. Mà kể cả tôi với Boji có xoàng xĩnh hơn, thì việc trụ vững trong một ``lớp chọn'' như 12A3 cũng chứng tỏ chúng tôi không đến mức hụt hơi.

``Ông biết gì chưa?'' Zetaro đột nhiên ghé sát lại, hạ thấp giọng đầy vẻ bí hiểm.

``\dots Gì?''

``Nghe đồn lớp mình sắp có thêm năm đứa mới chuyển vào đấy!''

Hakimi nhíu mày: ``Ngay lúc chúng ta sắp tốt nghiệp á?''

``\dots Hả?'' tôi giống như kẻ trên trời rơi xuống.

Khan cũng tặc lưỡi bất bình: ``Nghe vô lý vãi chưởng. Hôm nay là buổi học cuối cùng rồi còn gì.''

``Tao cũng thấy lạ. Nhưng mà thôi, kệ mịa nó đi,'' Hakimi xua tay.

``\dots Là sao vậy?'' Tôi vẫn chưa hiểu chuyện gì đang xảy ra.

Hakimi nhìn tôi bằng ánh mắt ái ngại: ``Vẫn còn ngái ngủ à? Hôm nay là buổi học cuối cùng của chúng ta trước khi nghỉ ôn thi đại học đấy.''

Không, cái này thì tôi hiểu. Vấn đề cốt lõi là\dots\ tại sao tôi lại bị ném quay về thời điểm này? Hơn nữa, trong ký ức của tôi, làm gì có chuyện lớp A3 đón thêm học sinh mới vào cái ngày sát nút này cơ chứ?

\dots Bây giờ tôi mới để ý kỹ. Tôi đang ngồi một mình một dãy. Cả bàn ngay phía trước và bàn thứ ba đằng sau tôi đều trống trơn. Không lẽ\dots\ có ai đó đã cố tình dàn xếp cho tôi ngồi ở cái vị trí kỳ quặc này?

Trong lúc tôi còn đang ngơ ngác giữa ranh giới thực giả, cả lớp bỗng nhiên nháo nhào chạy ùa về chỗ ngồi.

``Chết mịa. Mụ Quế vào!'' Zetaro hoảng hốt thì sầm.

``cô Katsura vào à?'' Tôi giật mình.

Khan vội vã thụt lùi về bàn dưới: ``Có gì tí xin cô ra rửa mặt nha ông Kuon.''

Mấy đứa bọn nó nhanh chóng an tọa ở các bàn phía dưới, chung tổ với tôi. Hakimi vẫn giữ nguyên nét mặt điềm nhiên như không, hướng ánh mắt về phía bục giảng.

``Được rồi, lớp ơi, trật tự nào!''

Một giáo viên nữ trung tuổi bước vào. Đó chính là cô Katsura, giáo viên chủ nhiệm kiêm giáo viên dạy Toán của lớp tôi. Cô nghiêm mặt nhìn quanh: ``Cô có điều muốn nói với cả lớp trước khi vào tiết học đây.''

Dù cô đã lên tiếng, vài đám học sinh ở cuối lớp vẫn tiếp tục xì xầm tán phét. Cô ngán ngẩm thở dài, rút chiếc thước gỗ to bản gõ mạnh xuống bàn giáo viên tạo ra âm thanh *BẠP! BẠP!* chói tai như một thói quen dường như đã đâm sâu từ thời dạy tiểu học.

``Trật tự nghe cô nói này! Hôm nay lớp chúng ta sẽ có năm bạn mới. Cô hy vọng cả lớp sẽ giúp đỡ các bạn ấy hết sức trong buổi học cuối cùng này.''

``Cô ơi!'' một học sinh ngồi dãy giữa giơ tay thắc mắc. ``Các bạn ấy chuyển vào lớp mình đúng lúc này thật ạ?''

Cô Katsura thở dài: ``Cái này thì cô cũng chịu. Chính các bạn ấy nằng nặc đòi học ở lớp mình đấy chứ.''

Zetaro ngay lập tức nhoài người lên, mắt sáng rực: ``Các bạn ấy là con trai hay con gái vậy cô?''

``Bất lịch sự quá ông tướng ơi\dots'' Một bạn nữ ngồi trên lườm cậu bạn cháy mặt.

``Vãi chưởng?'' Hakimi thốt lên.

``Thật sự luôn ấy.'' Tôi thở dài, cạn ngôn toàn tập.

Thế nhưng, trước câu hỏi có phần vô duyên của Zetaro, cô chủ nhiệm của chúng tôi lại trả lời một cách hết sức tự nhiên, không hề có chút bận tâm nào về sự thô lỗ của cậu bạn: ``Tất cả bọn họ đều là con gái.''

Chỉ cần câu nói đó như một lệnh khởi chạy, đám con trai trong lớp lập tức đồng loạt reo hò ầm ĩ, tay đập bàn chân đạp ghế rầm rầm như thể vừa trúng giải độc đắc vé số. Mấy bạn nữ ngồi phía trước cũng quay sang nhau ồ lên kinh ngạc, bàn tán rầm rầm về việc có người lạ đến lớp. Kẻ phấn khích nhất không ai khác chính là cái gã Zetaro ``đáng kính'' của chúng tôi, cậu ta thậm chí còn vỗ tay đôm đốm, mặt mày rạng rỡ như vừa được mở hội.

Hakimi, người vẫn ngồi yên vị ở dãy bàn bên cạnh, lập tức quay sang nhìn tôi. Chúng tôi trao đổi ánh mắt, trong đó lẫn lộn sự ngạc nhiên và hoài nghi.

``\hl{Tôi cảm thấy có gì đó sai sai\dots}'' tôi thì thầm nhỏ đủ để chỉ cậu ta nghe thấy, lòng đang dấy lên một cảm giác khó tả. Có gì đó không ổn chút nào: tại sao năm cô gái lại chuyển vào lớp đúng vào buổi học cuối cùng trước kỳ thi?

``\hl{Tôi cũng ngửi thấy mùi mờ ám rồi đấy\dots}'' Hakimi đáp lại, đôi mắt nheo lại như đang cố gắng phân tích tình hình, mép miệng nhếch lên một nửa đầy hoài nghi.

Lớp học lại bắt đầu loạn lên như cái chợ vỡ. Tiếng nói, tiếng cười, tiếng đập bàn hòa quyện vào nhau thành một mớ hỗn độn. Gần như lúc nào cũng vậy: chúng tôi là lớp 12A3, cái lớp được mệnh danh là ``nổi loạn nhất khối''. Chính vì thế mà trong các đợt xếp hạng kỷ luật hàng tuần của trường, lớp chúng tôi lúc nào cũng cũng bị xếp ở vị trí chót bảng (giỏi lắm thì ngoi lên được áp chót, nếu may mắn không có lớp nào gây chuyện lớn hơn). Các giáo viên bộ môn mỗi khi có tiết ở cái lớp ``chọn'' này đều phải tốn không biết bao nhiêu calo và hơi giọng chỉ để ổn định trật tự, có thầy cô còn phải mang theo cái micro cầm tay mới đủ sức át tiếng ồn ào của bọn tôi.

Thế nhưng kỳ lạ ở chỗ, dù kỷ luật thì tệ hại, về mặt học tập, điểm số các môn Toán, Lý, Anh của lớp tôi luôn chễm chệ trong top 3 toàn khối. Hầu hết mọi người trong lớp đều là những học sinh giỏi, chỉ là khi tụ tập lại thì ai cũng trút bỏ hết sự gò bó, trở thành những đứa trẻ nghịch ngợm nhất. Nhưng đó không phải là điều đáng bận tâm nhất lúc này; vấn đề là năm người lạ kia.

``Thôi, thôi! Cả lớp trật tự!'' cô Katsura phải đập thước gỗ mạnh mấy lần vào bàn giáo viên, tạo ra những tiếng bập bập chói tai để thu hút sự chú ý. ``Cô nói tiếp đây. Cả năm người họ\dots\ đều là những người rất nổi tiếng đấy.''

Câu nói này như một quả bom thứ hai ném vào cái nồi sôi đã sẵn. ``Vãi đạn!'' Boji, người vẫn luôn ngồi yên lặng ở cuối bàn, bỗng thốt lên không kiềm chế được.

``Thật hả cô!?'' một bạn nữ ở dãy giữa giơ tay lên, mặt đầy kinh ngạc.

``Chưa bao giờ nghĩ lớp mình lại có người nổi tiếng học chung đấy!'' một bạn khác cũng nối tiếng.

Đúng như dự đoán, cái vòm không gian này lại tiếp tục bùng nổ. Mọi người đều bàn tán, đoán già đoán non năm người đó là ai, làm gì mà nổi tiếng. Đám bạn tôi nháo nhào cả lên, có đứa còn hét lên hỏi liệu đó có phải là người nổi tiếng trên mạng không, có đứa lại đoán là ca sĩ, diễn viên. Lúc nào cũng thế, lớp tôi chẳng bao giờ chịu được một chút tò mò là lập tức ầm ĩ lên.

Thấy tình hình có vẻ càng lúc càng bất ổn, không có dấu hiệu dịu đi, tôi quyết định im lặng. Rút chiếc áo khoác đồng phục của mình ra, tôi chùm kín đầu, che hết cả mặt rồi gục xuống bàn, chẳng thèm để tâm gì tới khung cảnh hỗn loạn đang diễn ra trước mắt. Tôi muốn tránh khỏi tất cả sự ồn ào này, cố gắng suy nghĩ xem mình đang ở đâu, tại sao lại bị đưa quay lại thời điểm cấp ba, và năm người chuyển đến kia là ai.

Kijiru, cậu bạn vốn luôn tò mò về mọi thứ, nhô đầu ra khỏi cửa sổ sát hành lang, mắt hướng về phía cửa lớp: ``Vậy họ đâu rồi hả cô? Sao chưa vào?''

``Họ đang đứng chờ bên ngoài lớp đấy,'' cô Katsura đáp, vừa nói vừa hướng chiếc thước kẻ của mình ra phía cửa, ra hiệu cho những người đứng ngoài vào.

Một vài đứa ngồi dãy ngoài cùng sát hành lang, vốn đã chờ đợi từ lâu, vội kiễng chân lên, cố gắng ngó ra ngoài hành lang. Lát sau, một bạn gái reo lên phấn khích: ``Em thấy họ rồi cô ơi! Xinh lắm luôn ấy!''

Nghe câu đó, hội con trai trong lớp như bị thôi miên, lập tức nhào tới cửa sổ, đẩy nhau chen lấn để được nhìn thấy: ``Đâu? Đâu?'' ``Cho xem với!'' ``Tại sao mày che hết cửa sổ thế!'' Tiếng cười đùa, tiếng la hét lại tiếp tục vang lên, làm lớp học càng lúc càng ồn ào không khác gì một buổi hội chợ.

Riêng tôi thì vẫn cứ nằm úp mặt vào lớp áo khoác đồng phục, hai bàn tay bịt chặt lấy tai, gắng sức cách ly mình ra khỏi cơn hỗn loạn vô nghĩa đang lan tỏa khắp phòng học. Một cảm giác bồn chồn khó tả cứ thế dâng lên trong lồng ngực, như có một linh tính vô hình đang thì thầm rằng những người sắp bước vào đây sẽ không chỉ xáo trộn buổi học bình thường này, mà còn có thể lật tung cả cuộc đời tôi trong cơn mơ kỳ quặc này.

``T-Thôi được rồi, vào đi các con ơi!'' tiếng cô Katsura vang vọng từ bục giảng, hướng ra phía cửa lớp.

Và rồi, từng giọng nói nhỏ nhẹ nối nhau đáp lại từ bên ngoài hành lang.

``Vâng~~''

``Vâng!''

``Vâng\dots''

``Vâng ạ, peko!''

``Vào liền ạ!''

Những chất giọng đó\dots\ quá đỗi quen thuộc. Tôi đã nghe đi nghe lại chúng hàng trăm, hàng ngàn lần trong suốt thời gian làm việc ở \hll. Không có khả năng tôi nhầm lẫn được, dù chỉ là một chút.

Trái tim đập thình thịch như sắp bật ra khỏi lồng ngực. Tôi chậm rãi hé mắt từ dưới lớp áo khoác, không dám nhúc nhích mạnh, chỉ dám nhìn lén qua một khe nhỏ để xem những người vừa bước vào. Tôi vừa hồi hộp vừa lo lắng, chẳng thể ngờ giấc mơ phi lý này lại đang dần trở nên rõ ràng đến thế.

\lang{Nhật}

Dưới sự dẫn dắt của cô chủ nhiệm, năm người con gái đó lần lượt bước qua ngưỡng cửa lớp, nối tiếp nhau bước vào giữa không khí đang sôi nổi như nồi lẩu đang sôi của 12A3. Nổi bật đi đầu là Noel, Pekora và Marine - cả ba tiến vào với những bước chân đầy tự tin và dạn dĩ. Theo sát phía sau là Flare với phong thái từ tốn, dịu dàng. Cuối cùng, rụt rè và nép mình tận phía sau cùng là Rushia, vẻ mặt cô nàng hiện rõ sự e sợ và bối rối đến tội nghiệp.

Cả năm người họ đều khoác lên mình bộ đồng phục nữ của trường Anwa: chiếc áo phông trắng có cổ phối cùng chiếc đầm xếp ly màu xám trang nhã kéo dài tới trên đầu gối một chút. Ngoại trừ Cô đồng đang run run vì ngượng ngùng, bốn người còn lại đều nở nụ cười rạng rỡ chào cả lớp.

Sự xuất hiện của họ chẳng khác nào một mồi lửa ném vào thùng thuốc súng. Đám con trai trong lớp lập tức reo hò, huýt sáo vang trời. Thậm chí cái gã Zetaro còn phấn khích đến mức rời khỏi chỗ ngồi, đứng hẳn dậy vỗ tay bôm bốp.

``\textbf{Tôi biết bọn họ! Tôi biết bọn họ rồi!}'' cậu bạn Bách khoa ấy hét lên.

Kitaharu nhíu mày nhìn kỹ rồi quay sang xầm xì: ``Hình như\dots\ họ thuộc \hll\ đúng không nhỉ? Trong đó hình như là còn có oshi của mày trong đó đấy, Zentaro.''

Kijiru từ bàn bên nhếch môi cười mỉa: ``Khoản này thì chắc chắn có người trong nhóm mình hiểu rõ hơn tất cả chúng ta đấy.''

``Để tao gọi!'' Zetaro mắt sáng rực, quay phắt xuống dãy tôi ngồi, ``\textbf{Kuon ơi! Kuon!}''

Nghe thấy tên mình bị xướng lên giữa chốn đông người, tôi trùm áo khoác giả chết cũng phải khẽ giật mình một cái. Trong lòng tôi thầm gào thán: ``\textbf{\textit{Đừng có lôi tao vào cái vụ này chứ cái đồ đầu bình hỗn làm này\dots}}''

Hakimi ngồi ngay cạnh tôi tất nhiên không bỏ qua chi tiết đó. Cậu ta nheo mắt nhìn tôi đầy vẻ tò mò: ``Ông biết họ à?''

Tôi thở dài một tiếng qua lớp vải áo: ``\dots Biết. Biết rất rõ là đằng khác.''

``OK\dots? cậu bạn gật gù.

``Nhưng tôi không nói đâu,'' tôi củng cố lại tư thế nằm úp mặt xuống bàn. ``Để xem mọi người sẽ phản ứng thế nào.''

``\eng{\dots Được thôi, tùy ông thôi\dots}'' Hakimi đáp, giọng điệu thản nhiên như thường lệ.

Tất nhiên, tôi quyết định giấu tiệt chuyện mình chính là Tổng quản lý của bọn họ ở thế giới thực. Lý do rất đơn giản: nói ra vào lúc này thì có chó nó mới tin! Một thằng học sinh bình thường, suốt ngày lù khù ở bàn cuối mà lại là sếp của dàn idol nổi tiếng toàn cầu á? Nghe hư cấu chẳng khác nào mấy cuốn Light Novel rẻ tiền.

Nghĩ là làm, tôi tiếp tục áp dụng chiến thuật ``im lặng là vàng'', giả vờ như mình đã chìm sâu vào giấc ngủ.

Thế nhưng, đời không như là mơ. Lưới trời lồng lộng, tuy thưa mà khó thoát. Ngay khi cái gã Zetaro ăn hại kia réo tên tôi, cả lớp lập tức dồn toàn bộ sự chú ý về phía bàn của tôi.

``Ô-Ông Kuon ơi\dots''  một bạn học ngồi dãy bên cất tiếng hỏi.

``Bạn thuở nhỏ của ông đấy à!?'' một bạn học khác nhô đầu lên.

``Mày quen biết họ từ khi nào vậy hả Kuon?'' một ông con trai trong lớp ghen tị ra mặt.

``Thế mà lại giấu anh em\dots\ Ích kỷ vừa thôi nhé!'' một ông khác cay đắng bĩu môi.

Thấy tôi bị đưa vào thế gọng kìm, Hakimi thở dài một tiếng rồi lên tiếng giải vây: ``Thôi đi mấy ông ơi. Ông ấy không thoải mái với mấy cái bình luận vô căn cứ ấy đâu.''

Mặc cho dư luận vây quanh, tôi vẫn kiên định với bài ``im lặng là vàng'', coi như tất cả chỉ là tiếng gió thoảng qua tai. Cứ giả điếc như vậy cho chắc ăn. Chẳng ai làm gì được tôi cả.

Lúc này, Kijiru giơ tay hướng về phía bục giảng hỏi: ``Vậy các bạn là ai thế?''

Flare mỉm cười đặt tay lên ngực: ``Nhắc mới nhớ thì\dots''

``Chúng ta cũng phải giới thiệu bản thân đúng không ta?'' Noel tiếp lời.

Cô Katsura gật đầu hài lòng: ``Đúng rồi đó. Nào cả lớp, trật tự đi để cho các bạn ấy nói!''

Như có một phép thuật nào đó chi phối, toàn bộ cái lớp học vốn náo nhiệt như cái chợ vỡ bỗng nhiên im bặt hẳn. Dường như ai cũng tò mò muốn biết thân thế thực sự của năm cô gái đặc biệt này.

``Tất cả các bạn ở đây đều thuộc nhóm \hll\ Thế hệ thứ Ba\dots\ hoặc ít nhất, cô được bên đại diện giới thiệu như vậy,'' cô chủ nhiệm vừa giải thích vừa quay sang nhìn cả năm người, ``Được rồi, mấy đứa giới thiệu bản thân đi.''

Nhận được sự cho phép của giáo viên, các thành viên của \textit{\hll\ Fantasy} lần lượt tiến lên phía trước. Người đầu tiên xông xáo nhất không ai khác chính là cô nàng Thỏ ngọc. Pekora chống hai tay lên hông, đưa tay ra trước tạo dáng đầy năng lượng: ``Kon-peko, kon-peko, kon-peko! Thành viên thuộc thế hệ thứ ba của Hololive, mình là Usada Pekora-peko! Doumo, doumo!''

Ở góc bàn phía dưới, Hakimi gãi đầu ngơ ngác xầm xì với Boji: ``Sao tôi nghe cái từ cuối giống `almond-o' (hạt hạnh nhân) thế nhỉ?''

Boji gật gù tán thành: ``Ừ, tôi nghe cũng thấy giống hệt.''

``Pfft\dots\''tôi trùm áo khoác, không kiềm chế nổi mà bật cười phì một tiếng.

Chiếc tai thỏ của Pekora khẽ giật giật. Cô nàng ngay lập tức vô nhịp, cau mặt hướng về phía tôi mà đập bàn pằng pằng: ``Này nhá! Tôi nói là `doumo', chứ không phải là `almond-o' đâu đấy-peko!''

Tôi đành kéo chiếc áo khoác xuống một chút, lộ ra nửa khuôn mặt tà tưa: ``\dots Thì em nói cũng tựa tựa gần âm thế mà.''

``\textbf{Không phải cả anh cũng hùa theo họ chứ, Kuon-senpai!?}'' Pekora tức tối giậm chân.

``Senpai!?''

Ngay lập tức, một luồng điện xẹt qua toàn bộ phòng học. Cả lớp đồng loạt trố mắt nhìn tôi bằng ánh mắt kinh ngạc tột độ.

Trong lòng tôi thót lại một cái, mồ hôi hột bắt đầu tuôn ra như mưa: ``\textit{Vãi linh hồn! Bị bại lộ nhanh thế sao!?}''

Tôi đành thở dài, kéo hẳn cái mũ áo khoác xuống rồi ngẩng mặt lên, chạm ánh mắt với năm cô gái đang đứng trên bục giảng.

Kitaharu ngơ ngác quay sang hỏi tôi: `` `Senpai' ư? Họ là lớp dưới của ông à, Kuon?''

``Cũng\dots\ không hẳn.'' Tôi gãi đầu bối rối. ``Chuyện này nói ra thì dài lắm\dots''

Rồi tôi quay sang lườm cô nàng Thỏ ngọc: ``Mà sao em lại biết anh ở cái lớp này vậy, Peko-chan?''

Pekora nhếch môi cười đắc thắng, vểnh cằm lên: ``Bọn em đã biết anh học ở lớp này từ lâu rồi-peko!''

Marine đứng bên cạnh liền nghiêng người làm điệu, nháy mắt tinh nghịch: ``Thế nên bọn em mới chủ động đăng ký chọn đúng cái lớp này để học cùng anh đấy, Kuon-kun~!''

Lời thú nhận của Thuyền trưởng như một gáo nước sôi dội thẳng vào cái lớp A3 này. Hakimi quay sang nhìn tôi bằng ánh mắt hình viên đạn, sắc lẹm đến mức có thể cắt đứt kim loại: ``''...Thế là thế nào hả Kuon? Tôi yêu cầu ông một lời giải thích đàng hoàng đấy.''

Biết là chẳng thể nào giấu giếm thêm được nữa, tôi buông một hơi thở dài ngao quán. Thôi thì\dots\ đã biết rồi thì cho biết cho chót vậy.

Tôi dõng dạc lên tiếng, giọng điệu đầy vẻ chấp nhận số phận:

``Tôi đây chính là Tổng quản lý của bọn họ.''

Cả lớp - bao gồm cả cô giáo chủ nhiệm Katsura - đều há hốc miệng kinh ngạc. Các thành viên trên bục giảng thì gật gù liên tạch như chim gõ kiến để bảo chứng cho tuyên bố của tôi: ``Đúng rồi đó~! Kuon-senpai thật sự là Tổng quản lý của bọn mình!''

Zetaro thì rõ ràng đã biết điều này từ trước liền vỗ đùi đét một cái: ``Tôi biết ngay mà! Quả nhiên là thế!''

Boji thì ngơ ngác không tin vào mắt mình: ``Ô-Ông làm cái công việc đó từ khi nào vậy?''

``Được khoảng hai năm rồi,'' tôi uể oải đáp. ``Còn lý do tại sao ấy\dots\ nó phức tạp lắm. Hôm nào rảnh tôi sẽ nói sau.''

Kijiru nhíu mày thắc mắc: ``Nhưng mà\dots\ tại sao họ lại xuất hiện ở đây vào lúc này?''

``Tôi chịu!'' tôi vò đầu bứt tai, ``Đến chính tôi còn đếch hiểu tại sao tôi lại ngồi ở đây nữa là!''

``Là sao?'' Hakimi lại hỏi.

``\dots Vốn dĩ tôi là quản lý của họ ở thế giới thực. Thế rồi tối hôm qua đi ngủ, đến khi mở mắt ra thì chẳng biết tại sao lại lạc vào cái giấc mơ kỳ quặc này.''

Pekora áp hai tay lên má, đôi mắt lấp lánh như đang chìm đắm trong một bộ phim ngôn tình: ``Có thể\dots\ nó chính là `định mệnh' chăng, peko?''

``Không! Không đời nào nó lại vô lý và sến sẩm thế được\dots'' tôi xua tay gạt phắt đi, rồi lập tức đanh mặt lại, ``Với cả tổ lái vừa thôi cô nương.''

Cảm nhận được không khí trong lớp lại sắp vượt khỏi tầm kiểm soát, cô Katsura vội vàng giương chiếc thước gỗ lên, gõ mạnh xuống mặt bàn giáo viên với những tiếng động *BẠP! BẠP!* chói tai để kéo lại trật tự. Mặc dù vậy, ánh mắt cô vẫn không thể rời khỏi tôi, đong đầy sự ngỡ ngàng chưa dứt: ``Cô thực sự không ngờ rằng con lại có thể trưởng thành và giỏi giang đến thế, Kuon\dots''

Tôi vội xua tay, giọng nói lắp bắp vì sự bối rối bất ngờ: ``À không cô ơi! Chuyện này nó\dots nó cũng đến với con một cách cực kỳ đột ngột, con đâu có chuẩn bị gì đâu mà--''

Chưa kịp cho tôi nói hết câu thanh minh, đám bạn cùng bàn đã đồng loạt quay sang nhìn tôi với ánh mắt pha trộn giữa nghi ngờ và đố kỵ tột độ. Nhận thấy bầu không khí đang dần trở nên ngột ngạt, có thể bùng nổ bất cứ lúc nào, Hakimi liền nhanh trí chen ngang để đổi chủ đề, cứu tôi khỏi thế bí: ``Cái đó cứ để bàn sau đi. Các bạn mới còn lại cũng nên giới thiệu bản thân với cả lớp chứ nhỉ?''

Trong lòng tôi thầm thở phào nhẹ nhõm, như vừa trút được gánh nặng. Ơn giời có cậu bạn này ra tay giải vây, nếu không biết phải xử lý tình huống ra sao. Tôi thầm nhủ nếu có dịp gặp lại Hakimi ở Etsunan sau này, chắc chắn sẽ mời cậu ta ăn một bữa thịnh soạn để đáp ơn. ``Cảm ơn ông thật nhiều.''

Ngay sau lời đề nghị của Hakimi, cô nàng Hiệp sĩ đứng kế bên Pekora liền mỉm cười thân thiện, bước nhẹ nhàng lên phía trước bục giảng. Shirogane Noel đặt một tay lên ngực, nở ra nụ cười ấm áp như ánh nắng mùa xuân làm ấm lòng người. ``Konbanmassuru~ Mọi người gọi mình là Shirogane Noel, một thành viên của \hll\ Thế hệ thứ Ba nha~''

Ngồi dưới gốc bàn, tôi khoanh tay trước ngực, không nhịn được mà bày ra trò trêu chọc thường ngày, chêm vào một câu nhỏ đủ để mọi người xung quanh nghe thấy: ``\dots Và còn có biệt danh không thể nào nổi tiếng hơn: `Hiệp sĩ PON vô dụng'.''

Ngay khi cái biệt danh vừa bật ra khỏi miệng tôi, cô nàng Hiệp sĩ lập tức nhảy cẫng lên như vừa bị chạm vào dây thần kinh. Đôi má trắng hồng bỗng chốc đỏ ửng hết cả lên, toát lên vẻ tức giận khó che giấu. Noel giậm chân liên tục xuống sàn gạch, giọng nói run rẩy đầy vẻ oan ức: ``Em không hề PON cũng chẳng vô dụng đâu chứ! Sao mà anh cứ thích đi trêu em hoài vậy hả, Kuon-senpai!?''

Ngồi ở dãy bàn phía dưới, Zetaro nheo mắt nhìn chằm chằm vào Noel từ đầu đến chân, chẳng biết đang suy tính điều gì. Bỗng dưng, cậu bạn thốt lên một câu cực kỳ thành thật mà chẳng hề che giấu chút nào: ``Công nhận ngực bạn to thật đấy\dots''

Hakimi ngồi cạnh Zetaro lập tức câm nín, quay phắt sang nhìn cậu bạn với ánh mắt khó hiểu. Cậu ta khẽ thốt lên, giọng đầy ngạc nhiên: ``\dots Thật luôn hả ông?''

Bị ánh mắt của cả đám con trai trong lớp dồn vào, Noel chẳng còn biết giấu đi đâu nữa. Cô nàng vội vã nâng hai tay che chắn ngực lại, nhắm chặt cả mắt để tránh giao tiếp với bất kỳ ai. Noel kêu lên bằng giọng đầy thẹn thùng: ``Đ-Đừng có nhìn mình bằng ánh mắt như vậy chứ~!!''

Nhìn cảnh tượng vừa diễn ra, tôi chỉ còn cách lắc đầu thở dài, cố gắng cảnh báo đám bạn về sức mạnh thật sự của cô nàng Hiệp sĩ. Tôi vỗ nhẹ lên vai Hakimi, giọng nói đầy nghiêm trọng: ``Nói thế thôi chứ thực ra em ấy mạnh khủng khiếp đấy. Đến độ chỉ cần lỡ tay đấm nhẹ một cái là vỡ nát cả tường gạch luôn đấy.''

Nghe tôi nói vậy, Noel lập tức ngước mắt lên, đôi mắt còn ngân ngấn nước như sắp khóc. Cô nàng mếu máo, vội vàng giải thích với hy vọng cứu vãn lại hình ảnh của mình: ``Cái đó chỉ là em vô ý thôi mà!''

Lượt giới thiệu tiếp theo đến với thành viên còn lại của nhóm Thế hệ thứ Ba, và người đầu tiên bước lên bục giảng chính là Flare - cô nàng nửa yêu tinh với mái tóc vàng óng ả như những sợi nắng buổi chiều. Dáng đi của cô vừa thong thả vừa trang nhã, mang theo phong thái của một quý cô thực thụ, khác hẳn với sự năng động quá mức của cô bạn Thỏ ngọc vừa rồi.

Flare cúi đầu thật sâu, gửi đến cả lớp một cái chào hết mực chuẩn mực: ``Chào buổi chiều-nui~ (Kon-nui~), mình là Shiranui Flare, thuộc thế hệ thứ ba. Mong mọi người chiếu cố!''

Boji, vốn là người ít nói nhất trong đám bạn tôi, liền gật gù không ngớt: ``Bạn ấy trông có vẻ lễ phép đấy nhỉ.''

Kitaharu ngồi cạnh cậu bạn cũng lia mắt nhìn Flare không rời, thêm vào một câu đồng tình: ``Trông cũng có vẻ rất tốt bụng và dễ gần nữa.''

Cái gã Zetaro vốn hay phóng khoáng thì càng không kiềm được, cậu ta thốt lên bằng giọng khàn khàn sau khi nuốt một ngụm nước bọt to đùng: ``Đáng yêu thật đấy\dots''

Tôi cũng liền gật đầu liên tục, đồng lòng với tất cả những gì mấy đứa bạn vừa nói: ``Đúng. Tất cả những gì mấy ông nói đều đúng hết.''

Một cô bạn ngồi ở dãy giữa lớp bỗng giơ tay cao, mặt tràn đầy sự tò mò không che giấu: ``Bạn là yêu tinh đúng không?''

Nghe câu hỏi đột ngột, Flare hơi nhướng mày ngạc nhiên, đôi mắt long lanh chớp chớm vài lần như không hiểu sao người kia lại nhận ra thân phận của mình: ``Sao bạn lại biết hay vậy?''

Tôi chẳng đợi cô nàng kịp nói thêm gì liền nhếch môi lên tiếng giải đáp thay Flare, giọng trêu chọc không giấu giếm: ``Chắc là bởi đôi tai nhọn hoắt kia kìa. Elf nào mà chả như thế. Ngoại trừ việc em là nửa người, nửa yêu tinh.''

Nghe tôi trêu chọc, Flare cũng chẳng hề tức giận. Cô nàng chỉ quay sang nhìn tôi, nở ra một nụ cười nhẹ nhàng như gió thổi qua mặt hồ, toàn bộ vẻ mặt đầy sự dịu dàng khó tả.

Tiếp theo, cô nàng mái tóc đỏ rực như ngọn lửa, với chiếc bịt mắt cướp biển che nửa mắt trái, từ từ bước lên bục giảng. Houshou Marine cố tình hạ thấp giọng, tạo ra một âm vực trầm ấm mà đầy lôi cuốn, như muốn hút hồn bất kỳ ai lắng nghe.

``Ahoy! Mình là Thuyền trưởng Houshou Marine, thành viên của Hololive Thế hệ thứ Ba đây~!''

Cô nàng khoác tay tạo dáng đầy đường cong, rồi gửi một nụ hôn gió bay về phía cả lớp, hành động cởi mở khiến không khí phòng học lập tức sôi nổi hơn. Phần lớn đám con trai trong lớp lập tức ồ lên phấn khích, không ngừng huýt sáo thốt lên những lời khen ngợi lộ liễu.

Tuy nhiên, riêng nhóm ``Rách'' chúng tôi ngồi ở góc lớp vẫn giữ nguyên vẻ mặt thản nhiên, thậm chí còn nheo mắt nhìn cô nàng với sự hoài nghi rõ rệt. Vài học sinh nghiêm túc ở dãy bàn đầu cũng không có nhiều phản ứng, chỉ gật gù xem qua rồi quay sang bàn bạc nhỏ giọng với nhau.

Marine thấy phản ứng từ lớp không hề như những gì cô mong đợi, lập tức cau mày, mái tóc đỏ rực khẽ rung lên vì bất bình. Cô xị mặt xuống, đôi mắt nhìn thẳng về phía nhóm chúng tôi rồi lớn tiếng hỏi: ``Gì đây hả? Cái ánh mắt đó là sao?''

Tôi khoanh hai tay trước ngực, nhếch mép một nửa đầy trêu chọc, rồi lớn tiếng đáp lại câu hỏi của cô nàng: ``Em có chắc em là `thuyền trưởng' thật không vậy?''

``Chứ còn sao nữa!'' Marine chống hai tay lên hông, giọng nói cãi bướng không chịu thua, chỉ thẳng vào mặt tôi mà hỏi: ``Anh có ý kiến gì à!?''

Chưa kịp để tôi mở lời, Zetaro ngồi ngay dưới bục giảng đã phán một câu xanh rờn, lời nói thẳng thắn đến mức khiến cả lớp bật cười: ``Mình có cảm giác bạn giống như một bà già 70 tuổi đang cố cosplay làm thuyền trưởng thì đúng hơn á.''

``I-Im đi!'' Marine lập tức đỏ bừng cả mặt, hai tay chỉ trỏ lung tung vì hoảng hốt, giọng nói lắp bắp không ra hơi: ``Ông nghĩ cái kiểu gì vậy hả!?''

``Nó nói có sai quái đâu chứ.'' Tôi bật cười khoái chí, vỗ đùi liên tục vì phấn khích trước câu nói của cậu bạn Bách khoa. Sau đó tôi quay sang nhìn cậu bạn bằng ánh mắt thán phục, lớn tiếng khen: ``Mà Zetaro-senpai này, mày xem chừng cũng hiểu biết về mảng này lắm nhể?''

``Kuon-senpai\dots!'' Marine nổi giận đùng đùng, đôi mắt phừng phừng lửa, giậm chân mạnh xuống sàn lớp thô bạo, kêu lên: ``\textbf{Rút lại lời nói đó ngay cho em!}''

Tôi vờ như không nghe thấy lời đe dọa của cô nàng, tiếp tục ``bóc phốt'' những chuyện riêng của nàng Thuyền trưởng trước mặt toàn bộ lớp học để trêu chọc. Tôi vờ như đang suy nghĩ, rồi chậm rãi nói: ``À, và em ấy cũng từng tự tay làm một tựa game hồi còn học trung học nữa đấy. Game RPG gì đó thì phải\dots''

``!!!'' Senchou giật thót người, mặt cắt không còn một giọt máu, hai mắt tròn xoe kinh ngạc, vội vàng giơ tay ngăn cản tôi nói tiếp: ``C-Chờ đã--!''

Cả lớp lại một lần nữa trầm trồ nhìn Marine bằng ánh mắt ngưỡng mộ, những câu nói liên tục vang lên khắp phòng học, khiến không khí càng lúc càng sôi nổi. Một bạn học sinh ở dãy giữa hét lên: ``Vãi! Bạn ấy cũng biết làm game cơ à?''

Một bạn khác liền nối tiếng, mặt đầy ngưỡng mộ: ``Giỏi thế!''

Lại một bạn học sinh khác thốt lên, giọng đầy phấn khích: ``Đúng là không hối hận khi được học cái lớp này mà!''

Cuối cùng, một đứa con trai ngồi cạnh tôi ở tổ bên cạnh càu nhàu nhỏ, giọng đầy ghen tị: ``Thằng Kuon sở hữu dàn gái khủng thế! Thế mà không chia bớt cho anh em\dots''

Marine sau khi nghe hết những lời nói của cả lớp, mặt đỏ bừng từ tận cổ lên đến tận tai, hai tay vội vàng che lấy mặt vì ngượng ngùng. Cô nàng kêu lên nhỏ, giọng đầy oan ức: ``Đồ xấu tính! Sao lại đi nói cái điều đáng xấu hổ ấy ra cơ chứ\dots\ Sao anh ấy lại biết được chuyện đó hay vậy\dots''

Cô Katsura lắc đầu, nụ cười lan tỏa trên khuôn mặt như đang xem một màn trình diễn vui nhộn. Cô vẫy chiếc thước gỗ nhẹ nhàng để dẹp đi sự ồn ào còn sót lại sau màn trêu chọc Marine vừa rồi. ``Thôi, không trêu bạn ấy nữa các con. Còn bạn cuối cùng, giới thiệu cho cả lớp nghe đi nào.''

``V-Vâng ạ!'' Một tiếng kêu nhỏ xíu, rụt rè cất lên từ góc cuối hàng năm cô gái.

Giọng nói đó yếu ớt đến mức nếu lớp không vừa im lặng sau tiếng gõ của cô giáo, chắc chắn sẽ chẳng ai nghe thấy. Nó càng làm nổi bật sự khác biệt giữa cô nàng tóc xanh lá này với bốn người đồng đội đầy tự tin và tràn đầy năng lượng bên cạnh.

Khác hẳn với những người đi trước, Rushia tỏ ra vô cùng rụt rè, gần như muốn thu mình nhỏ lại trong chiếc áo đồng phục rộng. Chính xác hơn, cô đang hoảng loạn tột độ, đầu óc quay cuồng vì không biết phải diễn đạt thế nào để để lại ấn tượng tốt nhất với bạn mới.

``A-nô\dots\ Ê-tồ\dots'' nàng Pháp sư xoắn xít hai ngón tay vào nhau, móng tay không ngừng vân vê vạt áo đồng phục. Cặp mắt to tròn cứ cúi xuống nhìn sàn nhà, chẳng dám ngước lên nhìn ai trong lớp.

Marine đứng bên cạnh vỗ vai cổ vũ mạnh mẽ, cố gắng truyền cho cô nàng chút can đảm. ``Thôi nào Rushia. Bà chính là nhân vật chính của buổi học ngày hôm nay đấy.''

Noel cũng mỉm cười ấm áp, đôi mắt hiền lành hướng về Rushia như muốn nói rằng mọi thứ sẽ ổn thôi. ``Cố gắng lên nào, Rushia-chan!''

Tôi từ dưới bàn cuối lớp cất giọng ôn tồn, đủ để chỉ mình cô nàng nghe thấy. ``Cứ nói giống như những gì bản thân em vẫn hay nói hàng ngày là được mà.''

Rushia ngước mắt nhìn tôi, như thể vừa tìm thấy một tia hy vọng cuối cùng. Cô thở phào một cái rồi hắng giọng lấy lại sự tự tin bị đánh mất. ``Vâng! E-hèm\dots''

``Kon-rushi~ Mình là Uruha Rushia, thành viên Hololive Thế hệ thứ Ba. Rất hân hạnh được gặp gỡ mọi người!''

Cả lớp lập tức đáp lại cô bằng một tràng pháo tay ròn rã, to đến mức làm rung cả cửa sổ sát hành lang. Mọi tiếng ồn ào biến mất, thay vào đó là những ánh mắt ấm áp dành cho cô gái nhút nhát vừa đứng lên.

``Tốt rồi đấy,'' tôi gật đầu khen ngợi, thầm mừng vì cô nàng cuối cùng cũng dám mở lời. Không có gì phải lo lắng cả, mọi người đều đón nhận cô ấy nhiệt tình.

Zetaro ngồi dưới bục giảng thì thì thầm với cậu bạn bên cạnh, giọng đầy thiện cảm. ``Bạn ấy trông cũng lịch sự và ngoan ngoãn không kém gì Flare-chan ha.''

Bạn nữ ngồi bàn trên cậu ta cũng gật gù lia mắt nhìn Rushia không rời, miệng nở nụ cười dễ thương. ``Chưa kể giọng còn nhẹ nhàng nữa chứ! Trông dễ thương xỉu luôn ấy!''

Nghe những lời khen liên tục vang lên, nàng Pháp sư ngượng ngùng gãi gãi má, đôi má ửng hồng lên vì xấu hổ. ``C-Cảm ơn các bạn\dots\ Ngượng quá đi\dots''

Bạn nam ngồi dãy giữa hào hứng giơ tay lên, không che giấu sự tò mò cháy bỏng. ``Thế bạn tự giới thiệu công việc của mình như thế nào vậy?''

``À, m-mình là một cô đồng (Necromancer).'' Rushia ngập ngừng giải thích, tay vẫn không ngừng xoắn vắt áo. Cô ngước nhìn mọi người với ánh mắt lo lắng, sợ rằng họ sẽ sợ hãi trước năng lực đặc biệt của mình. ``Kiểu như\dots\ có thể triệu hồi và hồi sinh người chết ấy?''

Ngay lập tức, Kitaharu mắt sáng rực như vừa phát hiện ra kho báu, cậu không giấu được sự phấn khích. ``Thật á!? Bạn làm thử cho mọi người xem được không?''

Tôi giật mình, vội giơ tay ra hiệu muốn ngăn cản kịp thời. Tôi quá rõ điều gì sẽ xảy ra nếu để Rushia thực hiện phép thuật của mình trong lớp học này. ``Ờm, tôi không nghĩ là ông nên yêu cầu em ấy làm cái đó--''

``Được thôi!'' cô nàng tóc xanh tươi cười đáp lại ngay lập tức, làm tôi hoàn toàn sững sờ. Chẳng những không ngần ngại, cô còn tỏ ra vui mừng vì cuối cùng cũng có cơ hội thể hiện năng lực của mình.

Trước khi kịp có ai đưa ra lời ngăn cản nào, Rushia đã bắt đầu lẩm bẩm một câu thần chú cổ xưa, giọng nói nhỏ bé nhưng lại vang vọng một cách kỳ lạ trong toàn bộ căn phòng. Chỉ một giây sau, một cánh tay xương xẩu trắng bệch đột ngột nhô lên từ mặt sàn gạch, xuyên qua lớp gạch men cứng cáp như thể nó chỉ là đám bùn nhão mềm nhũn.

*RẮC! RẮC!*

Tiếng gạch nứt vỡ vang lên như một tiếng sét đánh giữa trời yên, và một bộ xương khô nguyên vẹn từ từ trồi lên hoàn toàn. Nó đứng sừng sững ngay trên bục giảng, thậm chí còn cao hơn cả cô Katsura đang đứng chết trân tại chỗ không thể tin được.

``\textbf{C-Cái quái gì thế này!?}'' Đám con trai trong lớp hoảng sợ nhảy bật dậy, có đứa còn làm lật cả ghế ngồi bởi sự kinh hãi tột độ. Tất cả ánh mắt đều đổ dồn vào bộ xương kỳ lạ kia, không ai dám tin vào những gì mắt đang thấy.

Rushia nghiêng đầu với vẻ mặt ngơ ngác, dường như hoàn toàn không hiểu tại sao mọi người lại phản ứng thái quá đến thế: ``Hửm? Chỉ là phép triệu hồi cơ bản của mình thôi mà. Mình tưởng bà đồng nào mà chẳng làm được thế này?''

Zetaro run cầm cập chỉ vào bộ xương, giọng nói lạc hẳn đi vì quá kinh hãi: ``L-Lần đầu tiên trong đời mình nhìn thấy một bà đồng thực thụ đấy\dots''

Tôi thở dài úp mặt vào lòng bàn tay, thầm hiểu rằng những rắc rối chỉ mới vừa bắt đầu. Tôi đã trải qua đủ chuyện với các cô nàng Thế hệ thứ Ba để biết rằng đây mới chỉ là món khai vị: ``Cái này vẫn còn bình thường chán\dots''

Kijiru tò mò hỏi thêm, tuy mặt có chút tái đi nhưng vẫn không giấu được sự thích thú lạ thường. ``Thế á? Vậy bạn còn làm được gì nữa không?''

``Mình có thể điều khiển cả một đám cùng một lúc luôn đấy!'' Rushia đáp với vẻ mặt đầy tự hào, ngực nở ra khi nhận được sự chú ý của cả lớp. Cô nàng hoàn toàn quên đi sự nhút nhát ban đầu, giờ đây chỉ còn niềm vui khi được thể hiện khả năng đặc biệt của mình.

Ngay khi câu nói vừa dứt, cô nàng giơ cao bàn tay lên. *RẮC! RẮC! RẮC!* Hàng chục bộ xương khô khác đồng loạt chui lên từ mặt sàn, lấp kín toàn bộ khu vực bục giảng. Cả lớp kinh ngạc đến mức không thể thốt nên lời.

Rushia vẫy tay ra lệnh với giọng đầy phấn khích: ``Nào các Fandead, nhảy một điệu múa cho mọi người xem đi!''

Đám bộ xương lập tức tuân lệnh. Cùng lúc đó, không biết từ đâu vang lên một bản nhạc EDM nồng nhiệt, sùng sục vang vọng khắp phòng học. Đám xương kia bắt đầu nhảy những động tác quẩy sung cực kỳ điêu luyện, khiến ai cũng phải tròn mắt ngạc nhiên.

Bất giác, một vài cậu con trai trong lớp cũng lao lên bục giảng, nhiệt tình quẩy cùng đám bộ xương. Bên dưới, những người khác cũng bắt đầu vỗ tay hô vang theo nhịp điệu. Chẳng mấy chốc, phòng học 12A3 vốn trang nghiêm đã biến thành một sàn nhảy ngầm chính hiệu.

Hakimi ngồi cạnh tôi trố mắt nhìn cảnh tượng trước mắt, hoàn toàn câm nín không nói nên lời.: ``\dots Oắt-đờ-heo\dots''

Tôi chỉ còn cách lắc đầu ngao ngán trước những gì đang diễn ra: ``Đó cũng là một kỹ năng đặc biệt của em ấy nữa đấy\dots''

Cô Katsura trên bục giảng hoàn toàn bất lực. Cô chỉ có thể đập chiếc thước gỗ liên hồi với những tiếng *BẠP! BẠP!* để cố gắng kéo lại trật tự.

``\dots Được rồi, và đó là Uruha Rushia! Cảm ơn con! Cả lớp về chỗ ngồi ngay cho cô! Chúng ta sẽ làm quen thêm ở tiết sinh hoạt! Còn bây giờ thì để cô xếp chỗ cho các con\dots''

Tôi tặc lưỡi phán trước với thảm trạng sắp xảy ra: ``Dám cá là họ sẽ ngồi ở các vị trí xung quanh mình cho xem, lạ gì nữa.''

Quả nhiên không ngoài dự đoán, giáo viên chủ nhiệm chỉ tay về phía khu vực bàn của tôi: ``Các con xuống ngồi ở những bàn trống xung quanh vị trí của Hikawa Kuon nhé.''

``Vâng~~!'' cả năm cô gái đáp lại một cách dõng dạc, đầy phấn khích.

Trong khi đó, toàn bộ đám con trai trong lớp (trừ nhóm ``Rách'' chúng tôi) đều lườm tôi bằng ánh mắt hình viên đạn, đầy sự ngỡ ngàng và ghen tị cháy da cháy thịt.

Noel hào hứng chạy xuống từ bục giảng: ``Wai~ Tôi muốn ngồi cạnh Flare!''

Flare dịch vào trong, sẵn sàng đón bạn mình: ``Ờm, được chứ. Bà cứ việc.''

Marine thì lau mồ hôi trên trán, áp sát vào Rushia với vẻ mặt nũng nịu: ``Tôi muốn ngồi chung với Rushia cơ~ Có được không vậy?''

Rushia đổ mồ hôi hột, vội vàng xua tay né tránh sự tiếp cận của Thuyền trưởng: ``Bà tởm quá đấy Marine! Sao bà không sang ngồi với Pekora ấy?''

``Nhưng mà tôi không muốn ngồi với thỏ đâu~ Đi mà~? Xin bà đấy~?'' Senchou chu môi làm nũng, không chịu bỏ cuộc.

``Kinh quá đi mất! Thôi được rồi, tôi ngồi với bà là được chứ gì!?'' nàng Pháp sư thở dài bất lực trước sự dai dẳng của bạn mình: ``Thật là\dots\ Ít nhất còn dễ thở hơn là ngồi cạnh cái tên hay gây sự kia.''

Pekora đứng ngay cạnh bàn tôi, cau có xị mặt xuống với vẻ không hài lòng: ``\dots Này, vậy là giờ em phải ngồi cạnh anh đấy à-peko?''

Tôi ngước nhìn cô nàng Thỏ ngọc, thừa nhận tình hình hiện tại: ``Dường như là thế. Trông em có vẻ không thích nhỉ?''

``Không phải thế!'' cô ấy giậm chân thanh minh, sợ tôi hiểu lầm: ``Chỉ là\dots\ em chẳng hiểu tại sao bọn họ lại xếp chỗ như thế này nữa-peko\dots''

Tôi ngửa đầu ra sau ghế, cũng cảm thấy bất lực trước sự sắp xếp kỳ lạ này: ``Bắc thang lên hỏi ông trời đi em ạ, chứ anh là anh bó toàn bộ não rồi.''

Hakimi ghé sát tai tôi thì thầm, đầy tò mò về mối quan hệ của các cô gái: ``\dots Hỏi không liên quan nhé Kuon, nhưng mà bọn họ\dots\ có cặp với nhau thật à?''

Tôi quay sang nhìn Hakimi với ánh mắt nghiêm túc, cảnh báo bạn mình về những gì sắp xảy ra: ``Ở \hll\ thì chuyện `teetei' (cơm chó) này là cơm bữa rồi. Nếu ông không muốn ăn phát ngấy thì tôi khuyên thật lòng là nên quay đầu sang hướng khác đi. Nghiêm túc đấy.''

Kijiru ngồi dưới xuýt xoa khi nghe đến chuyện ``ăn cơm chó''.

``Ăn nhiều nên quen rồi bạn ạ,'' tôi thở dài cam chịu với số phận của mình: ``Đừng lo cho tôi về cái khoản đó. Lo cho tôi về cái đám con trai đang nhìn tôi như muốn ăn tươi nuốt sống kia kìa.''

Tôi vừa dứt câu, sát khí từ đám con trai trong lớp lại bùng lên dữ dội hơn nữa. Từng ánh mắt hừng hực lửa ghen tức dồn thẳng vào tôi như muốn thiêu đốt.

Thế nhưng, chưa kịp để tụi nó làm gì tôi, cô Katsura đã đập mạnh thước kẻ xuống bàn để ngăn chặn hỗn loạn: ``Thôi ngay đi! Kuon ơi! Cả mấy đứa kia nữa, cứ sồn sồn hết cả lên thế!'' cô nhắc nhở với tâm trạng chẳng mấy dễ chịu: ``Buổi học cuối cùng rồi mà mấy đứa cứ nghịch như cái chợ vỡ thế hả?''

Trái ngược với sự chịu đựng của tôi, năm thành viên Thế hệ thứ Ba lại đồng loạt quay xuống, tỏa ra thứ không khí bảo vệ tôi một cách cực kỳ thái quá (mặc dù tôi chẳng hề yêu cầu hay cần thiết).

Pekora đập bàn đứng dậy, tỏa ra sự quyết tâm bảo vệ tôi: ``Đừng có lo, Kuon-senpai! Bọn em sẽ bảo vệ anh-peko!''

Noel giơ nắm đấm siết chặt, sẵn sàng đối đầu với bất kỳ ai: ``Em có thể đấm gãy xương mấy đứa kia nếu anh muốn!''

Marine cũng giơ nắm đấm ra, không chịu thua kém ai: ``Marine đây cũng không chịu thua đâu nha~!''

Sợ nhất là cô đồng Rushia. Mái tóc xanh của cô nàng khẽ rủ xuống, đôi mắt lóe lên luồng sát khí đen ngòm kinh hoàng: ``\dots Kẻ nào dám làm hại đến anh ấy\dots\ ta sẽ cho xiên thẳng tay.''

Cả bốn người bọn họ - đặc biệt là từ Cô đồng gọi hồn - tỏa ra thứ áp lực đè nặng khiến toàn bộ đám con trai trong lớp lập tức im bặt, nuốt nước bọt ừng ực vì khiếp sợ trước sự đáng sợ đó.

Hakimi đổ mồ hôi hột, thì thầm vào tai tôi với giọng run rẩy: ``\dots Đám con gái của ông đáng sợ quá đấy, Kuon ơi\dots''

Boji lắc đầu ngao quán, đồng tình với cảm nhận của Hakimi: ``Trông họ có vẻ cũng không phải dạng thích đùa đâu\dots''

Koheki ngồi bàn cuối lúc này mới cất giọng lạnh lùng, nhận thấy sự nguy hiểm của các cô gái: ``Tốt nhất là đừng có đứa nào dại dột đụng vào thằng Kuon là được. Tao thì tao chẳng có ý kiến gì về họ cả.''

Zetaro ngồi kế bên liền cà khịa Koheki, trêu chọc bạn mình về chuyện tình cảm: ``Xem cái thằng vừa bị bồ đá nói kìa~''

Cậu bạn Hóa học lập tức nổi gân cổ, xoay người lại định dạy dỗ Zetaro: ``Tao đấm bỏ mẹ mày bây giờ! Thằng đầu bình này!''

Flare vội vàng can ngăn cãi vã, cố gắng giữ hòa khí trong lớp: ``Thôi nào mọi người\dots\ Đánh nhau không hay ho gì đâu mà.''

Koheki hắng giọng, quay đi chấp nhận lời can ngăn của cô nàng yêu tinh: ``Nhỏ yêu tinh đó nói đúng đấy. Tao cũng chẳng thừa hơi mà muốn tẩn nó một trận đâu.''

Tôi bất lực chắp tay vái các cô nương, cầu xin các cô bình tĩnh lại: ``Anh xin mấy đứa đấy. Cứ ngồi trật tự và ngoan ngoãn giùm anh cái. Anh không muốn cái lớp này tan đàn xẻ nghé chỉ vì mấy đứa đâu.''

Cô Katsura hắng giọng một tiếng lớn để lấy lại uy quyền, vội vàng đổi chủ đề để rời khỏi tình thế khó xử. ``V-Vậy là ta đã ổn định xong chỗ ngồi rồi ha? Bây giờ cả lớp lấy đề ôn tập ra, chúng ta sẽ tiếp tục chữa từ đề số 12.''

Cô nhìn sang năm học sinh mới vừa chuyển vào, nhắc nhở về việc học tập: ``À, mấy đứa vừa vào lớp có thể xem chung đề với các bạn xung quanh nhé.''

Chẳng cần suy nghĩ đến nửa giây, cả năm người bọn họ đồng loạt đứng dậy, nhào tới vây kín xung quanh chiếc bàn của tôi để đòi xem chung đề. Họ dường như chẳng hề quan tâm đến những bạn khác xung quanh.

Đám con trai trong lớp đứng nhìn cảnh tượng đó mà máu ghen tức sôi lên sùng sục. Riêng tôi thì chỉ biết buông tay thở dài, chấp nhận số phận của một kẻ thủ khoa ngành gánh nạn:

``\dots Mấy đứa muốn làm gì thì làm đi.''

\subsubsection*{\phantomsection\label{ep-C7.2}}
\addcontentsline{toc}{subsubsection}{{\sen{-Partie 2-}} {\caviar Những tiết học cuối cùng}}
\stpt{-Partie 2-}\stcap{Những tiết học cuối cùng}

Vì bằng một cách thần kỳ nào đó mà tôi lại (dường như) bị giật lùi thời gian quay về đúng cái giai đoạn ôn thi đại học khốc liệt, nên ngay sau khi cả lũ vừa hoàn thành xong kỳ thi học kỳ, cả trường đã lập tức bước vào chiến dịch ``cày ải'' các môn trọng tâm. Đối với năm cô nàng Thế hệ thứ Ba gồm Pekora, Noel, Flare, Marine và Rushia, đây chắc chắn là một trải nghiệm có một không hai khi lần đầu tiên trong đời được nếm thử không khí ngột ngạt của một ``lò ôn thi'' đích thực tại xứ Etsunan này.

Nói qua một chút về cái trường Anwa này, tất cả các khối trong trường được phân chia rõ ràng thành hai ban: Ban A và Ban D.

Lớp A3 của chúng tôi thuộc Ban A - cái ban chuyên cày các môn khoa học tự nhiên như Toán, Vật lý, Hóa học và Sinh học. Tất nhiên là vẫn không thể thiếu hai môn bắt buộc là Ngữ văn và Tiếng Anh. Nhìn qua cái danh sách này thì có thể thấy ngay, mấy môn tự nhiên nặng đô này hoàn toàn không phải là sở trường của năm cô idol nhà \hll.

Trong khi đó, Ban D ở dãy đối diện lại chuyên về các môn khoa học xã hội như Văn, Lịch sử, Địa lý và Giáo dục công dân. Cả đống môn ngập tràn chữ nghĩa mà đám học sinh chúng tôi vẫn hay gọi vui bằng cái tên thân thương nhưng mỉa mai là ban ``học thuộc''.

Một ngày học của chúng tôi kéo dài nửa ngày với tổng cộng năm tiết. Buổi học hôm nay là buổi cuối cùng trước khi nghỉ cuối tuần, cộng thêm việc tiết cuối cùng là tiết Sinh hoạt lớp, nên tôi bấm ngón tay nhẩm tính\dots\ hôm nay chắc chắn là thứ Sáu.

Thứ Sáu hôm đó có năm tiết nối tiếp nhau: Toán, Lý, Hóa, Tiếng Anh và Sinh hoạt.

Khai thật thì trong các môn tự nhiên, môn Hóa học chính là ``kẻ thù truyền kiếp'' của tôi, tôi học dở nhất môn này. Dở nhì chính là môn Vật lý. Riêng với môn Toán, phần Đại số tôi cân hết, làm ``mượt như sunsilk'', nhưng cứ hễ đụng đến Hình học không gian là não tôi lại bật chế độ ``ngắt cầu giao''. Còn môn Tiếng Anh thì\dots\ khỏi phải bàn, cứ nhìn cái cách tôi chém gió bắn Tiếng Anh rần rật như gió mùa đông bắc với đối tác và các thành viên ở văn phòng \hll\ hàng ngày là đủ hiểu trình độ của tôi rồi.

Nói về tiết Toán của cô Katsura, cách lớp A3 chúng tôi học phải gọi là độc nhất vô nhị. Quy trình trả lời câu hỏi ôn tập không hề bình thường như các lớp khác. Để trả lời các câu hỏi trong đề ôn, mỗi người chúng tôi phải móc điện thoại ra, gõ đáp án rồi gửi thẳng vào nhóm trò chuyện Messenger có cái tên cực kỳ học thức: ``A3 có học.''

Cái thú vị (và cũng là thảm họa) của cái nhóm này nằm ở chỗ: nếu có ai đó ngứa tay chọn một đáp án lệch quẻ so với số đông còn lại, cô giáo chủ nhiệm sẽ lập tức réo tên nạn nhân đó đứng dậy giải thích tại sao mình lại chọn như vậy. Mà đến 80-90\% các trường hợp ``một mình một ngựa'' như thế đều là khoanh sai bét, thành ra người đó sẽ được hưởng trọn một tràng cười nghiêng ngả từ cả lớp, tất nhiên là cười đùa vui vẻ chứ không hề có ý ác ý gì.

Tôi chợt nhớ ra năm cô nàng đang ngồi quanh mình làm gì đã có trong cái nhóm chat bựa này. Tôi giơ tay lên cất tiếng: ``Con thưa cô!''

Cô Katsura ngẩng đầu lên khỏi cuốn giáo án: ``Sao thế con, Kuon?''

Tôi chỉ tay về phía năm cô gái đang ngồi xung quanh mình: ``Các bạn ấy chưa được thêm vào nhóm chat của lớp thì trả lời kiểu gì ạ?''

Cô chủ nhiệm chớp mắt nhìn năm học sinh mới: ``Mấy đứa không dùng Mes (Messenger) sao?''

Pekora đưa tay gãi đầu, đôi tai thỏ khẽ giật giật đáp: ``Dạ\dots\ Bọn con ở bên đó chủ yếu dùng Disc*rd với L*NE để nhắn tin với làm việc là chính thôi ạ-peko\dots''

Flare cũng mỉm cười ngại ngùng gật đầu: ``Thật ra bọn em còn chưa bao giờ nghe đến cái ứng dụng đó nữa ạ.''

Cô Katsura trầm ngâm một chút rồi phẩy tay: ``Hừm\dots\ Vậy thì các con dùng chung điện thoại với Kuon nhé.''

Cả năm cô nàng nghe xong liền đồng loạt gật đầu đồng ý tắp lừ. Nhưng Pekora lại nghiêng đầu thắc mắc: ``Nhưng mà nhắn kiểu gì cho cô biết là ai chọn đáp án nào vậy ạ-peko?''

Cô Katsura mỉm cười hướng dẫn: ``Cứ gõ kiểu: `Rushia: A', `Marine: D' rồi gửi vào nhóm là được.'' Chỉ đơn giản như vậy thôi.

Thề, cái tiết Toán này lắm pha oái oăm không chịu nổi. Mới đến nửa buổi mà đã có tình huống dở khóc dở cười xảy ra.

``Câu số 16, các con chọn phương án gì nào?'' giọng cô chủ nhiệm vang lên.

Màn hình điện thoại của tôi lập tức nhảy tin nhắn liên tục. Một hàng dài chữ ``B'' xuất hiện rần rật trên nhóm chat của lớp. Tôi cũng nhanh tay gõ ``Kuon: B''. Bốn cô nàng Pekora, Flare, Marine, Rushia nhìn đáp án của tôi rồi cũng bấm chọn B theo.

Duy chỉ có một cái dòng tin nhắn vô cùng lạc quẻ hiện lên: ``Noel: A''.

``Shirogane-san! Tại sao con lại chọn đáp án A vậy?'' quả nhiên, cô Katsura lập tức soi ra ngay cái đáp án đi ngược lại với cả vũ trụ kia.

Nghe thấy tên mình bị cô giáo réo lên giữa lớp, cô nàng Hiệp sĩ giật bắn mình, luống cuống đứng bật dậy. Mặt em ấy đỏ bừng, ấp úng trả lời với sự thật thà đến ngây thơ: ``Dạ\dots\ dạ thưa cô! Tại\dots\ tại em thấy trong bốn đáp án, cái số nào to nhất thì em chọn thôi ạ!''

*PHỤT! HA HA HA!*

Cả lớp A3 lập tức bùng nổ trong một tràng cười ngặt nghẹo. Đám con trai vỗ bàn bôm bốp, còn nàng Hiệp sĩ thì ngơ ngác đứng chết lặng tại chỗ, đôi mắt chớp chớp không hiểu rốt cuộc mình đã nói sai điều gì.

Pekora lắc đầu thở dài, khoanh tay phán một câu đầy vẻ bề trên: ``Tự tin gớm nhỉ-peko. Nhưng mà sai lè ra rồi đấy nhé!''

Hakimi ngồi cạnh tôi cũng phải giật giật khóe mắt, thì thầm thán phục: ``\dots Wow. Đến cái phép tính cơ bản của học sinh cấp 3 mà cũng không làm được luôn à?''

Tôi nhếch môi cười tà tưa, ghé sát tai cậu bạn ngồi cạnh thì thầm tiết lộ bí mật quốc gia: ``\hl{Ông có muốn biết bài thi Toán đợt kiểm tra hồi Tết năm ngoái em ấy được bao nhiêu điểm không?}''

Hakimi ngơ ngác quay sang nhìn tôi: ``\hl{\dots Thật á? Ông còn tổ chức kiểm tra ấy luôn à?}''

``\hl{Thật 100\%. Đoán thử xem được mấy điểm?}''

Noel đứng phía trên nghe thấy tôi manh nha bóc phốt liền giật thót người, quay xuống chắp hai tay mếu máo van xin: ``K-Kuon-senpai! Đừng có nói chuyện đó ra trước mặt mọi người mà~!''

Hakimi gãi cằm suy đoán: ``\hl{\dots 4 điểm à?}''

Tôi lắc đầu, đưa hai ngón tay lên lắc lắc: ``\hl{Tệ hơn thế nhiều. 28. Trên thang điểm 100. Tính ra chưa nổi 3 điểm hệ 10 nữa.}''

Cậu bạn câm nín toàn tập, cất giọng bái phục: ``\hl{Vãi chưởng\dots{} Đúng chuẩn PON (vụng về) không trượt phát nào.}''

``Đ-Đừng có bóc phốt em như thế mà Kuon-senpai!'' Noel úp mặt xuống bàn khóc không thành tiếng.

Tất nhiên, bên cạnh những màn diễn hài đỉnh cao của cô nàng Hiệp sĩ, cũng có những khoảnh khắc khiến cả cái lớp A3 này phải há hốc miệng kinh ngạc.

``Câu 48, đáp án là gì nào các con?'' cô Katsura cất giọng hỏi.

Lần này, nhóm chat ``A3 có học'' im ắng một cách lạ kỳ. Chỉ có rải rác vài ba tin nhắn nhảy lên. Đây là một câu hỏi thuộc phần Hình học không gian nâng cao, đúng ngay cái ``vùng tối'' trong bộ não của tôi. Nhìn cái hình chóp chồng chéo góc và khoảng cách mà tôi muốn sang chấn tâm lý, hoàn toàn không biết phải bắt đầu từ đâu.

Trong khi tôi đang vò đầu bứt tai định bấm chọn lụi, thì Flare ngồi bên cạnh khẽ chòm sang, ngón tay thon dài chỉ vào màn hình của tôi gõ nhẹ: ``Flare: C''.

Một lúc sau, cô Katsura gật đầu hài lòng công bố: ``Đáp án đúng cho câu 48 này là C! Trong lớp mình có bạn nào chọn C không?''

Flare nhẹ nhàng giơ tay lên, cất giọng trong trẻo: ``Dạ, em ạ.''

Toàn bộ ánh mắt trong lớp lập tức dồn phắt về phía cô nàng Bán yêu tinh. Đám con trai trố mắt nhìn Flare với vẻ ngưỡng mộ không giấu giếm. Boji trầm ồ lên: ``Kinh thế nhở!''

Kitaharu cũng lắc đầu thán phục: ``Vãi chưởng thật chứ, câu đó khó vãi ra đấy!''

Nàng Yêu tinh tóc vàng mỉm cười dịu dàng, khiêm tốn đáp lại những ánh mắt bái phục chung quanh: ``Có gì đâu mọi người ơi\dots\ Bài này chỉ cần xoay góc tư duy một chút là ra ngay đáp án rồi mà.''

Tôi quay sang nhìn em ấy, nuốt nước bọt cái ừng ực: ``Bài đó là dạng bài lấy điểm 9-10 cực khó đấy em ơi\dots''

Flare chớp mắt ngơ ngác nhìn tôi: ``Ể? Thật ạ?''

Tôi thở dài: ``Anh là anh chuẩn bị quay random (ngẫu nhiên) rồi đấy. Nhưng mà nếu anh khoanh lụi thì chắc anh cũng chọn C thôi.''

Pekora nghiêng đầu tò mò: ``Thế á-peko?''

``Ừ. Một ông thầy dạy Toán online nổi tiếng mà anh xem hồi trước có bảo, nếu đi thi gặp câu hình không biết làm thì cứ đè toàn bộ đáp án C mà khoanh.''

Rushia ồ lên một tiếng: ``Hể~\dots\ Ra là vậy sao\dots''

Cô Katsura trên bục giảng nghe vậy liền hào hứng mời: ``Vậy thì Shiranui-san, con có thể lên bảng giải chi tiết câu này cho cả lớp cùng xem được không?''

``Vâng ạ!'' nàng Elf tóc vàng đáp lời rất nhanh rồi hăng hái bước lên bục giảng.

Cầm viên phấn trắng trên tay, cô nàng tháo vát vẽ lại hình rồi từng bước triển khai các công thức logic một cách cực kỳ mượt mà. Trước sự ngỡ ngàng của toàn thể học sinh ban A, em ấy giải xong bài toán hái điểm 10 đó mà chẳng tốn quá nhiều thời gian.

Khi Flare bước trở về chỗ ngồi, tôi giơ một ngón tay cái lên khen ngợi em ấy. Nàng Elf mỉm cười e ấp, dịu dàng đáp lại tôi với sự nhường nhịn: ``Em không giỏi đến thế đâu mà\dots\ Chẳng qua là may mắn thôi.''

Tất nhiên, màn tương tác thân thiết đó của tôi và Flare lại tiếp tục hứng trọn những ánh mắt hừng hực lửa ghen tị từ đám con trai trong lớp. Tôi thở dài thầm nghĩ trong đầu: ``\textit{Haiz\dots\ Ngày hôm nay chắc chắn sẽ là một ngày dài dằng dặc đây...}''

\ct

Tiếp sau tiết Toán chính là giờ Vật lý, một trong những tiết học mà tôi luôn cảm thấy e ngại nhất trong toàn bộ thời gian cấp ba của mình.

Thành thật mà nói, ba môn Vật lý, Hóa học và Sinh học trong kỳ thi tốt nghiệp THPT đều được gộp chung vào một bài thi Khoa học Tự nhiên tổng hợp, với đúng 40 câu hỏi cho từng môn. Những câu hỏi ở mức cơ bản thì tôi vẫn có thể xử lý trôi chảy, nhưng cứ gặp phải những dạng bài nâng cao, tôi thường phải đọc đi đọc lại mãi mới có thể hiểu được đề bài muốn nói gì.

Dường như những khó khăn mà tôi luôn gặp phải chẳng hề tồn tại với Pekora.

Cánh tay nhỏ nhắn của nàng Thỏ cứ liên tục vươn lên giành quyền trả lời từng câu hỏi trong bộ đề ôn tập, không để ai kịp đuổi kịp. Tôi thật sự không ngờ rằng cô nàng lại giỏi môn Lý đến vậy, thậm chí tốc độ phân tích và đưa ra đáp án còn vượt trội hơn cả hai cậu bạn học khá nhất lớp là Hakimi và Khan.

Giáo viên bộ môn vừa mới mở miệng đọc câu hỏi, giọng nói còn chưa vơi hết: ``Câu 21\dots''

``B ạ, peko!'' Pekora đáp liền chẳng cần mất một giây suy nghĩ.

Cậu bạn Bác sĩ (tương lai) ngồi phía trước há hốc mồm không thể khép lại, giọng còn chưa hết sốc: ``Nhanh vãi!''

Hakimi quay phắt sang nhìn tôi, đôi mắt mở to hết cỡ đến mức có thể lăn ra khỏi quầng mắt: ``Thật sự luôn á Kuon? Sao bạn ấy lại biết nhanh đến thế''

``Đừng có nhìn tôi,'' tôi vừa nói vừa vội vàng xua tay. ``Tôi cũng mới phát hiện ra sự bất thường này được mấy phút thôi.''

``Câu 22--'' tiếng đọc đề của cô giáo vừa cất lên.

``C ạ!'' lần này Pekora lại cướp lời nhanh hơn cả.

Tôi nuốt nước bọt ừng ực, chỉ biết thầm cầu cứu trong lòng: ``Chậm thôi em ơi, để anh còn kịp đọc đề mà tiêu hóa nữa chứ\dots''

Cô giáo dạy Lý bất ngờ khựng lại vài giây, sau đó thở dài một tiếng rõ to: ``Câu 23 này\dots\ Usada-san để cho bạn khác trong lớp trả lời được không em?''

Pekora lập tức xỉu ngang trên bàn, mặt xịu xuống trông thật đáng thương: ``Ể~?''

Đúng vào lúc ấy, bộ đề ôn tập lại rơi vào một câu hỏi vận dụng cao nhất. Cả lớp lập tức chìm vào không gian im lặng đến mức nghe được cả tiếng kim rơi, chẳng ai dám ho he hay cọ quậy gây ồn ào.

Cô giáo đảo mắt nhìn lướt qua từng khu vực trong lớp, sau đó mới cất tiếng hỏi lại một cách ôn tồn: ``Có bạn nào biết làm câu này không? Ngoài bạn Usada-san ra?''

Sự im lặng nặng nề vẫn là câu trả lời duy nhất mà cô giáo nhận được.

Cô nhẹ nhàng thở dài, như thể đã đoán trước được tình hình: ``Vậy thì Usada-san, em lên trả lời cho cả lớp cùng nghe nhé.''

Ngay lập tức, nàng Thỏ tóc xanh nước biển lập tức hồi phục năng lượng như được sạc pin, giọng nói hồ hởi vang lên: ``C ạ!''

Giáo viên bộ môn chỉ có thể đáp lại bằng chất giọng pha chút mệt mỏi: ``Chính xác.''

Và cứ thế, chỉ một mình Pekora đã ``cân'' sạch ít nhất ba phần tư toàn bộ bộ đề ôn tập của tiết học hôm đó. Cả lớp chẳng ai nói nên lời, chỉ biết nhìn nhau trố mắt trước màn trình diễn chưa từng có của cô nàng, duy nhất chỉ có những tiếng thì thầm nhỏ lẻ vang lên: ``Bạn ấy ghê vãi chưởng\dots''

Bản thân tôi cũng bất ngờ đến mức bật ngửa khỏi ghế. Nhưng rõ ràng, những điều khiến tôi phải kinh ngạc trong ngày hôm nay vẫn chưa dừng lại ở đó...

\ct

Rồi cũng đến giờ Hóa học, môn học đầy rẫy những công thức và phản ứng hiểm hóc mà bất kỳ học sinh khối tự nhiên nào cũng phải phát sợ, và dĩ nhiên, tôi cũng không ngoại lệ. Lần này, nhân vật chính xuất hiện trên sàn diễn lại là Rushia.

Cô giáo bộ môn vừa mở miệng đọc đề: ``Chất nào sau đây không phải là axit? Đáp án A-''

``C ạ!'' nàng Pháp sư cất giọng nhẹ nhàng nhưng đầy sự chắc chắn, không chút do dự.

``Chính xác. Câu tiếp theo\dots''

Koheki, người ngồi ở bàn cuối cùng của lớp, thì thầm với bạn bên cạnh: ``\dots Nhanh vãi.''

Kihataru thì ôm đầu lắc lư không thể tin được: ``Não không load kịp luôn ấy\dots\ Bạn ấy nhảy đáp án nhanh quá!''

Một bạn nữ ngồi gần đó cũng trố mắt khen ngợi không giấu giếm: ``Rushia-chan giỏi thế!''

Rushia cười ngượng ngùng, hai má hơi ửng hồng lên vì được khen: ``Có gì đâu ạ, mình ăn may thôi mà.''

Trong khi mọi người còn đang kinh ngạc thì tôi đã gục mặt xuống bàn, chuẩn bị chìm vào giấc ngủ ngắn, thi thoảng lắm mới ngóc đầu dậy để xem tiết học đã kết thúc chưa.

Pekora ngồi bên cạnh thấy vậy liền chọc chọc vào lưng tôi, trêu ghẹo: ``Khó quá nên tắt nguồn luôn rồi hả, peko?''

Không chỉ những người đã chú ý từ đầu, mà cả cậu bạn suốt ngày chỉ biết cắm mặt vào game trên điện thoại cũng phải ngẩng đầu lên. Cậu ấy cũng không thể giấu được sự kinh ngạc trước những câu trả lời vừa nhanh lẹ lại chính xác tuyệt đối của Rushia.

Thầy giáo Hóa tiếp tục giọng đều đều: ``Tiếp theo, câu 31. Có bao nhiêu phát biểu đúng?''

Rushia lập tức giơ tay lên, không để ai kịp chần chừ: ``Em ạ!''

``Uruha-san, em trả lời đi.''

``Ba ạ. Là các ý (I), (III) và (V).''

Vị thầy giáo gật đầu hài lòng lắm: ``Chính xác. Em giải thích ngắn gọn cho cả lớp cùng nghe xem nào.''

Koheki ngồi dưới đứng hình vài giây, sau đó bật cười cay đắng: ``\dots Vãi chưởng thật.''

Noel thì mắt sáng rỡ như thấy báu vật, quay sang lay nhẹ áo Rushia: ``Rushia-chan, tí giảng lại cho mình khúc này với!''

Khan, cậu bạn học giỏi nhất lớp về các môn tự nhiên, cũng phải gật gù công nhận: ``Rushia tưởng bánh bèo mà học môn này siêu đấy chứ.''

Còn tôi? Tôi vẫn bận ``giao tiếp'' với Chu Tông trong giấc mơ của mình, tiếng ngáy nhỏ vẫn vang lên.

Tôi mơ màng nghĩ trong đầu, chẳng lẽ vì em ấy là một bà đồng chuyên gọi hồn nên mới học giỏi môn Hóa học này đến thế? Những con số và công thức hóa học kia chẳng khác nào những câu thần chú bí ẩn gì đâu. Dù gì thì đây cũng là môn sở đoản của tôi, nên tôi hoàn toàn không có ý kiến gì về màn thể hiện bá đạo này cả.

\ct

Tiết tiếng Anh bước đến ngay sau đó, môn học mà tôi và Hakimi tự tin nhất. Thế nhưng niềm vui ấy không hề chia sẻ được với mọi người trong lớp, khi tiếng Anh luôn là một trong những rào cản lớn nhất mà bất kỳ học sinh nào cũng phải đối mặt.

Những thành viên của thế hệ thứ Ba (hay cả\hll\ nói chung) thì càng gặp khó khăn hơn cả, với môn này gần như là cơn ác mộng không bao giờ kết thúc. Tôi biết rõ trình độ tiếng Anh trung bình của người Nhật không quá cao, nhưng chứng kiến họ vật lộn từng chữ, tôi vừa thấy thương hại lại vừa không nhịn được cười.

Một ví dụ điển hình xảy ra khi cả lớp đang cùng nhau giải các câu hỏi trong đề ôn tập.

Giáo viên tiếng Anh đặt câu hỏi: ``Câu số 25, các em chọn đáp án nào?''

Cả lớp đồng thanh hô lên: ``D ạ!''

``Rất tốt. Tiếp theo, câu 26, cô mời một bạn\dots''

Thỉnh thoảng, cô giáo sẽ gọi ngẫu nhiên một học sinh đứng lên trả lời và giải thích lý do chọn đáp án của mình. Chính vì ``trò chơi xổ số'' này mà nhiều lúc tôi, Hakimi, hay cả cô con gái của giáo viên chủ nhiệm - đúng, nhỏ đó cũng ở trong cái lớp 12A3 này - phải đứng lên gánh vác thay cả lớp.

Phải thừa nhận đây lại là một mô-típ anime quen thuộc: con gái của giáo viên học cùng lớp với nhân vật chính. Nhưng cô bạn ấy chẳng phải là nhân vật quan trọng mà chúng ta cần để ý đến lúc này.

À, tôi đang nói đến đâu rồi nhỉ? Quay lại chuyện chính nào.

Giáo viên tiếng Anh lại cất giọng: ``\dots Houshou Marine-san, em thử trả lời câu này nhé?''

Marine giật bắn người, mặt lúng túng đứng dậy vội vàng: ``V-Vâng! Đáp án là `ê-rích-sừn' (erection) ạ!''

Tôi phản xạ tự nhiên sửa phát âm ngay lập tức: ``Là `/ɪˈrek.ʃən/' mới đúng. Trọng âm nằm ở âm tiết thứ hai.'' Cô giáo lập tức gõ thước xuống bàn: ``Kuon, để yên cho bạn trả lời đi. Houshou-san, đọc lại theo cô nào.''

Tôi có một thói quen khá xấu là cứ thấy ai phát âm sai tiếng Anh là miệng lại tự động ``nói leo'' sửa lỗi, dù rõ ràng hành động đó chẳng lịch sự chút nào.

Hakimi ghé sát tai tôi thì thầm nhỏ: ``\hl{Này, đọc theo phiên âm Katakana của bà ấy cũng tạm được mà.}''

``Xin lỗi,'' tôi cười gãi đầu, ``\hl{tôi không kiềm được, cứ nghe ai đọc sai là tự động sửa. Cái tính cảnh sát chính tả nó vào người rồi.}''

Noel ngồi bên cạnh tròn mắt ngạc nhiên: ``\hl{Giờ em mới biết anh có cái tính này đấy\dots}''

Giáo viên tiếng Anh gật đầu hài lòng: ``Đúng rồi đấy Houshou-san. Em ngồi xuống đi, câu tiếp theo\dots''

Marine quay sang nhìn tôi với ánh mắt tủi thân, mếu máo: ``\dots Anh biết em không giỏi tiếng Anh rồi mà. Sao anh cứ bắt bẻ em làm em thấy khó chịu quá.''

Tôi vội vàng phân trần để cô nàng không hiểu lầm: ``Ở \hll\ gần như mọi người đều như thế cả thôi mà. Anh chỉ vô tình nói thật thôi.''

Pekora phồng má chọc ghẹo tôi, không vừa lòng: ``Ít nhất cũng đừng bêu xấu bọn em trước cả lớp chứ, peko. Anh có nhiều pha khiến bọn em mất mặt trước cả lớp rồi đấy.''

Rushia ngồi cạnh chỉ biết im lặng, gật đầu đồng tình với bạn mình, ái ngại nhìn tôi.

Chúng tôi tiếp tục giải thêm khoảng mười câu hỏi nữa. Luôn luôn là tôi hay cô bạn con gái giáo viên kia đứng lên gánh team trả lời. Mọi chuyện vẫn diễn ra bình thường cho đến khi\dots

Giáo viên lật sang trang mới của đề bài: ``Đến phần đọc hiểu thứ hai rồi\dots\ Houshou Marine-san, em tiếp tục trả lời các câu từ đây nhé.''

Nàng Thuyền trưởng giật mình hoảng hốt, không ngờ mình lại bị gọi lần nữa: ``Ể? E-Em nữa ạ?''

Tôi quay sang nhắc nhở cô nàng, vỗ vai an ủi: ``Chứ sao nữa. Câu từ 36 đến 40 đấy. Lần này anh hứa sẽ không nói gì cả, yên tâm.''

Senchou hít một hơi sâu để lấy lại bình tĩnh, gật đầu mạnh: ``R-Rồi. *e hèm* Câu 36, em chọn B ạ.''

Cô giáo gật đầu, yêu cầu thêm: ``Em đọc cả câu chứa đáp án đó lên cho cả lớp cùng nghe nhé.''

Marine bắt đầu ngập ngừng, mồ hôi hột bắt đầu đổ ra vì áp lực: ``À, a-nô, ê-tồ\dots `Hi oan tu se hít nâu-lếch tu e-vờ-ri-oăn' (He wants to share his knowledge to everyone)\dots''

Tôi phải dùng tay bụm chặt miệng để cố nhịn cười, không làm phiền cô nàng.

Hakimi ngồi cạnh lắc đầu ngán ngẩm, nhắc nhở tôi nhỏ: ``\hl{\dots Ít nhất tôi vẫn hiểu bà ấy đang đọc gì. Đừng cười nữa Kuon.}''

Tôi chẳng phải là người duy nhất khó nhịn cười. Một vài bạn trong lớp cũng bắt đầu tủm tỉm cười sau mỗi từ tiếng Anh lắp bắp của Senchou, với chất giọng Katakana đặc trưng. Ngay cả cô giáo trên bục cũng phải quay mặt đi để giấu tiếng cười.

Marine vẫn kiên trì đọc tiếp phần còn lại của bài: ``\dots Từ `rokku' (rock) trong câu trên đồng nghĩa với `shi-nê' (shine) ạ.''

Hakimi giật thót người, không tin vào tai mình: ``\dots Hả? Cái gì cơ?''

Boji ngồi phía trên cũng quay xuống tròn mắt, nhắc lại câu nói của Marine: ``Bạn vừa bảo người ta đi chết đi (shine) đấy à?''

Marine hoảng hốt xua tay giải thích, mặt đỏ bừng: ``K-Không phải! Mình đang định đọc cái từ `shi-nê' này nè, nhưng mà\dots''

Cả tôi và Hakimi cùng đồng thanh sửa phát âm cho cô nàng: ``\textbf{Là `/ʃaɪn/'!}''

Nàng Thuyền trưởng ngớ người ra một lúc rồi mới gật gù hiểu ra: ``À, ừ rồi, `Sai-nờ'!''

Tôi bất lực thở dài, cố gắng động viên: ``\dots Ờ thì, ít ra lần này nó cũng gần đúng rồi đấy.''

Cả lớp lại bùng nổ trong tiếng cười nghiêng ngả. Senchou đỏ mặt xấu hổ đến chín cả người, em ấy gục thụp xuống bàn, hai tay ôm chặt lấy đầu rên rẩm: ``Mọi người đừng có trêu mình nữa mà\dots''

Noel ngồi cạnh thở dài tự an ủi, thấy mình may mắn hơn: ``Có vẻ như mình không phải là người duy nhất gặp khó khăn với môn này rồi.''

Marine vẫn mếu máo van xin mọi người: ``Xin mọi người đấy, đừng chọc tôi nữa\dots''

Flare quay sang vỗ vỗ lưng an ủi bạn: ``Bà ổn chứ Marine?''

``Nhìn thế này mà bà nghĩ tôi ổn được hả? Xấu hổ muốn độn thổ luôn rồi đây này\dots'' Thuyền trưởng rên rỉ, vẫn không dừng ngẩng đầu lên.

Tôi thấy hơi áy náy nên lên tiếng an ủi cô nàng, hy vọng cô ấy bớt xấu hổ: ``Chuyện đó bình thường mà em. Nhiều lúc anh còn làm mấy pha `đi vào lòng đất' hơn thế nhiều.''

Hakimi lập tức bóc phốt tôi ngay trước mặt mọi người: ``Kiểu mấy lần ông phán `bài này dễ lắm, chỉ cần áp công thức là ra' trong khi bài tập khó vãi cả sít ra đúng không?''

Tôi cười trừ gãi cổ, không dám phủ nhận lời bạn nói: ``Kiểu kiểu thế đấy\dots''

Marine vẫn gục mặt xuống bàn, giọng lý nhí đáp lời tôi: ``Nói thế cũng chẳng làm em cảm thấy khá hơn đâu\dots''

Thuyền trưởng cứ thế úp mặt xuống bàn cho đến khi tiếng chuông báo hết tiết vang lên, mãi không dám ngẩng đầu lên.

Vậy là cả năm thành viên của thế hệ thứ ba Hololive đều đã bộc lộ hết trình độ học vấn của mình trước mặt cả lớp. Flare, Rushia và Pekora thì xuất sắc đến không tưởng, còn Noel và Marine lại đóng vai trò ``chốt hạ'' với những pha hài hước làm cả lớp cười đau bụng. Dĩ nhiên mọi người chỉ cười vui vẻ chứ chẳng hề có ý xấu, đây cũng là buổi học cuối cùng của thời học sinh mà.

Và giờ đây, tiết học mà tất cả chúng tôi mong chờ nhất đã đến: giờ Sinh hoạt lớp cuối cùng trong đời học sinh - ít nhất là theo những gì tôi biết ở thời điểm hiện tại.

\subsubsection*{\phantomsection\label{ep-C7.3}}
\addcontentsline{toc}{subsubsection}{{\sen{-Partie 3-}} {\caviar Cuộc chiến ngắn ngủi}}
\stpt{-Partie 3-}\stcap{Cuộc chiến ngắn ngủi}

Giờ giải lao cuối cùng trước tiết Sinh hoạt lớp - buổi học cuối cùng trong quãng đời cấp ba của tôi - cũng chính thức bắt đầu.

Sau mỗi hai tiết học kéo dài liền mạch, trường Anwa luôn dành ra mười lăm phút nghỉ ngơi để chúng tôi thả lỏng cơ thể sau những giờ tập trung. Những tiếng cười nói ríu rít lập tức lan tỏa khắp các góc lớp, nhóm kéo nhau xuống căng-tin xếp hàng mua đồ ăn vặt, nhóm khác xách giày chạy thục mạng xuống sân sau để tranh thủ đá vài quả bóng trước khi hết giờ. Không khí trong lớp chợt trở nên náo nhiệt hệt như một khu chợ nhỏ.

Còn tôi ư? Bình thường tôi chẳng hề tham gia vào đám đông ồn ào đó. Tôi luôn tìm một góc khuất yên tĩnh ở cuối lớp, lôi điện thoại ra bấm lướt một mình để tận hưởng khoảng thời gian riêng quý giá này. Đối với tôi, mười lăm phút yên ả ấy chính là thiên đường nhỏ bé mà chỉ riêng tôi mới có quyền sở hữu.

Nhưng hôm nay, mọi kế hoạch đều đảo lộn hoàn toàn. Thay vì thong thả lướt báo mạng hay mở game ra giết thời gian như mọi ngày, tôi lại phải ngồi gãi đầu gãi tai, xoay sở đáp lại đủ thứ câu chuyện lặt vặt từ năm cô gái Thế hệ thứ Ba đang vây kín quanh bàn của tôi. Cảm giác ấy thật sự chẳng khác nào một quản lý đang họp thăm hỏi tình hình nhân sự giữa giờ làm việc bận rộn.

Tôi cố gắng bắt nhịp theo dòng chuyện phiếm vô đầu của các cô nàng, chẳng mấy chốc chủ đề câu chuyện đã dần xoay quành về khoảng thời gian họ còn đi học ở quê nhà. Trời đất ơi, tôi còn chẳng bao giờ ngờ được những người bạn cùng lớp bất ngờ này lại từng có một thời thanh xuân đi học ly kỳ và khác biệt đến thế.

Pekora gật gù lia lịa, đôi tai thỏ trên đầu khẽ rung rung theo từng cử động: ``Công nhận giờ học của bọn anh khác xa với bọn em thật đó, peko.''

Tôi không khỏi tò mò, đặt câu hỏi mà lăn tăn trong đầu từ nãy: ``Thế hồi đó các em ôn thi như thế nào?''

Nàng Thỏ nhún vai một cái vô tư, chẳng có gì phải bận tâm: ``Nơi em sống làm gì có thi cử đâu, peko. Cứ học xong chương trình là ra trường thôi.''

Noel cũng gãi cằm suy nghĩ hồi lâu, cố gắng nhớ lại chi tiết về kỳ thi thời xưa: ``Bọn em thì chỉ có kỳ thi sát hạch hiệp sĩ thôi. Chủ yếu là kiểm tra thể lực với dùng tay chân là chính.''

Rushia khẽ chớp mắt vài lần, giọng nhỏ nhẹ như đang sờ lại những kỷ niệm cũ: ``Em thì phải vượt qua bài thi thực hành phép thuật\dots\ Nhưng mà lâu lắm rồi, chắc phải gần nghìn năm trước nên em cũng chẳng nhớ rõ nữa.''

Flare mỉm cười dịu dàng quay sang phía tôi, ánh mắt trông có vẻ nhớ về thời còn đi học: ``Em thì chắc giống anh nhất. Cũng đến trường, giải đủ các dạng đề rồi thi đại học. Ít nhất là ở Nhật Bản thì quy trình nó là như vậy.''

``Ra là vậy, mỗi đứa một kiểu\dots'' tôi gật gù lắng nghe, rồi đứng dậy vươn vai kéo giãn cơ thể. ``Xin lỗi mấy đứa nhé, anh ra ngoài rửa mặt cái cho tỉnh táo.''

Nàng Hiệp sĩ vui vẻ đáp lời, tay vẫn đang vẽ linh tinh vào quyển vở: ``Vâng, anh cứ đi đi ạ.''

Nàng Yêu tinh vàng cũng không bỏ qua cơ hội trêu ghẹo, chỉ vào gương mặt còn ngái ngủ của tôi: ``Thấy sau tiết Hóa vừa rồi trông anh vẫn còn ngái ngủ lắm đấy.''

``Ừ\dots'' tôi thở dài cười trừ, chẳng thể nào phủ nhận được. ``Anh xin phép nhé.''

Tôi bước nhanh ra khỏi lớp, đi thẳng về phía nhà vệ sinh ở cuối hành lang. Tôi tạt chút nước lạnh lên mặt, cảm giác mát lạnh lập tức xua tan đi cơn buồn ngủ còn vương vấn, cái đầu bớt đi cảm giác mông mông mây mây. Dù có ý định ``giải quyết nỗi buồn'' bằng cách ngủ thêm một chút khi quay về lớp, nhưng thấy chưa đến mức cấp bách, tôi đành thong thả túc tắc đi bộ quay trở lại.

\dots Dĩ nhiên, tuổi học trò làm sao mà thiếu được những xung đột bất ngờ, tình huống dở khóc dở cười mà tôi sắp gặp phải chính là một ví dụ điển hình. Suốt mấy năm ngồi trên ghế nhà trường, tôi cũng chưa từng gặp phải chuyện tương tự, bởi trường Anwa nổi tiếng với kỷ luật siết chặt đến mức đáng sợ. Ngay cả lớp A3 quậy phá nhất trường này cũng chẳng dám đốt cháy cả khu học xá như các thành viên thường làm ở công ty.

Vừa bước chân ra khỏi nhà vệ sinh chưa được mấy bước, một giọng nam trầm khàn khàn thình lình vang lên chặn đường tôi: ``\textbf{Ê, thằng kia!}''

Tôi ngẩng mặt lên đối mặt với kẻ vừa gọi.

Trước mặt tôi là một gã đầu trọc với làn da ngăm đen, gương mặt bặm trợn đầy mụn, tay cầm lăm lăm cây thước gỗ to đùng mà giáo viên hay dùng để kẻ bảng. Phía sau hắn là bốn tên tay sai đô con khác, vây quanh tạo thành một bức tường người chặn kín lối đi.

Tôi chẳng hề quen biết bất kỳ ai trong số bọn chúng. Nhìn phong cách ăn mặc và thái độ ngang ngược, có lẽ là học sinh lớp chuyên về thể thao ở lớp bên cạnh. Đám này rõ ràng thích giải quyết mọi chuyện bằng nắm đấm, một kỹ năng mà tôi tự nhận mình xếp cuối cả trường. Vì vậy, tôi quyết định chọn cách hòa giải nhẹ nhàng nhất.

Tôi thản nhiên đáp lời, giọng không chút nao núng: ``Phải, thì sao?''

Ngay lập tức, đám côn đồ cùng nhau bật cười hô hố, tiếng cười đầy miệt thị như thể vừa nghe thấy một câu chuyện cười hay nhất đời: ``Một thằng bé nhỏ tí như mày mà cũng đòi bám đuôi năm đứa con gái đẹp kia á!? Đừng có mơ!''

Tôi chỉ có thể thở dài ngao ngán, chẳng thể hiểu nổi đầu óc gã này đang chứa đựng những hoang tưởng gì mà tự cao đến thế.

Gã đầu trọc vung vẩy cây thước gỗ trước mặt tôi, giọng đầy vẻ khoe khoang: ``Nghe cho kỹ vào! Tao mới là người đủ đẳng cấp để tán tỉnh mấy em chân dài đó! Mày thì xứng đáng cái gì?''

``Thật ra tôi thậm chí còn chẳng hẹn hò với ai trong số họ cả,'' tôi bình tĩnh giải thích, hy vọng gã ấy sẽ tỉnh táo lại.

``\textbf{Tao đếch cần biết!}'' hắn gầm lên, thước gỗ đập thình thịch vào sàn nhà. ``Khôn hồn thì biến xa mấy đứa đó ra, nghe rõ chưa!?''

Đám tay sai phía sau được đà lập tức phụ họa theo, tiếng la hét vang dội cả hành lang: ``Đúng đó!'', ``Biến đi thằng ranh!'', ``Cút xéo khỏi đường trước khi đại ca tao nổi điên!''

Thấy bọn chúng càng lúc càng quá quắt, tôi không kiềm chế được mà buột miệng: ``Nói thật đấy\dots\ Mấy ông tỉnh lại giùm cái đi. Tôi nghĩ họ sẽ chẳng bao giờ thèm nói chuyện với mấy kẻ vũ phu, chỉ biết dọa nạt người yếu thế như các ông đâu.''

Gã đầu trọc lập tức nhướng mày, hất hàm đầy vẻ thách thức: ``Mày vừa nói cái quái gì cơ? Lại nói lần nữa xem nào?''

Tôi biết thừa hắn đang cố tình kiếm chuyện gây sự, nhưng tôi vẫn lặp lại câu vừa rồi bằng giọng điệu mỉa mai không chút che giấu: ``Bị điếc à? Tôi bảo là, họ sẽ chẳng bao giờ buồn tiếp chuyện với những kẻ hung hăng, thích giơ thước dọa người như các ông đâu.''

Lời vừa dứt, những đường gân xanh trên trán gã kia đã nổi lên cuồn cuộn, trông giống như sắp bật ra khỏi da.

``Nếu các ông định đánh tôi thì bỏ cuộc đi. Tôi chịu đau giỏi hơn các ông tưởng nhiều đấy.'' Tôi thản nhiên nói, trong đầu bất giác nhớ lại những lần bị các cô gái trêu chọc đến mức muốn đập đầu vào tường, cả về thể chất lẫn tinh thần.

Gã cầm đầu xốc lại cây thước gỗ, cười lớn cùng đám đàn em: ``Mày á!? Có mỗi một mình mày bé tẹo như con mắm thế này thì làm được cái trò trống gì!? Đừng có tự cao tự đại!''

``\dots Chắc chưa?'' tôi điềm tĩnh hỏi ngược lại, không hề bị động trước những lời đe dọa.

Hắn chỉ thẳng ngón tay vào mặt tôi, giọng đầy vẻ khinh bỉ: ``Mày nhìn lại mày xem! Xung quanh mày có ai đi theo bảo vệ không? Mày dựa vào cái quái gì mà dám mạnh miệng thế hả!?''

``Anh ấy nói đúng đấy.'' Một giọng nữ trong trẻo nhưng đầy quyền lực bỗng cất lên từ phía sau lưng bọn chúng.

``Chúng mày nên suy nghĩ lại trước khi quá muộn thì hơn.'' Một giọng nam khác cũng bồi thêm, vững chãi và chắc nịch.

\dots Đúng là một mô-típ quen thuộc không thể thiếu trong mọi câu chuyện. Khi nhân vật chính bị dồn vào đường cùng, kiểu gì cũng sẽ có các ``đấng cứu thế'' xuất hiện đúng lúc. Tôi không hề hoảng sợ từ nãy đến giờ chính là vì đã đoán trước được kịch bản này sẽ diễn ra.

Những người bước ra từ phía sau đám côn đồ không ai khác, chính là năm cô gái của Thế hệ thứ Ba, đi cùng với Hakimi và Zetaro -hai cậu bạn thân nhất của tôi trong lớp.

Hakimi bước lên trước, đút hai tay vào túi quần, ánh mắt sắc lẹm như dao găm: ``Mày thử động vào một cọng tóc của nó xem? Tao cho mày một phát đi chầu ông bà luôn đấy.''

Pekora cũng phồng má khẩy mũi, giọng đầy vẻ bất mãn: ``Bọn tao có bao giờ nói là sẽ hẹn hò với mày đâu mà đòi hỏi, peko!''

\dots Xem chừng một cuộc đối đầu lớn sắp nổ ra ở hành lang trường đây rồi.

Tên cầm đầu thấy dàn cô gái xuất hiện thì lập tức đổi thái độ, giọng nói dở cợt nhả dở tán tỉnh: ``Ái chà chà, các cô em ở đây miệng lưỡi cũng gắt gớm nhỉ? Hay là đi theo anh xuống căng-tin uống chút trà đàm đạo nhé~?''

Marine lập tức bước tới, hất mái tóc nâu đỏ đầy vẻ tự tin: ``Câu đó phải để ta nói mới đúng chứ! Đừng có mơ màng hão huyền!''

Tôi đứng ở đằng sau tròn mắt ngơ ngác, không ngờ cô nàng lại dũng cảm đến thế: ``\dots Thật luôn?''

Tên đầu trọc lập tức khoái chí cười khẩy, tưởng rằng cô nàng thực sự muốn đi theo mình: ``Thấy chưa? Chính cô ta muốn đi theo tao--''

``Nhưng phải là đi `chà bông' (đánh đập) cho nhừ xương mới đúng nhé!'' Marine cắt ngang lời hắn đầy vẻ bực bội.

Vừa dứt lời, Thuyền trưởng đã lao tới tát một cú nảy đom dóm mắt vào má tên tay sai đứng gần nhất, khiến cả đám côn đồ tròn mắt kinh ngạc. Chắc là ở công ty thường xuyên bị Rushia tát để trừng phạt những trò đùa dởm, nên giờ em ấy tranh thủ xả hết stress lên người khác rồi.

Tên cầm đầu nhanh chóng lấy lại bình tĩnh, giơ cây thước gỗ lên đầy vẻ đe dọa: ``Cô em trông ngon nghẻ thế mà tay chân cũng bạo lực phết nhỉ. Được thôi, nếu các cô đã muốn chơi kiểu này!''

Hắn quát lớn ra lệnh cho đám đàn em, giọng đầy vẻ thị uy: ``\textbf{Lên cho tao! Tất cả cùng xông vào đánh bọn chúng!}''

Nếu bạn đang tưởng tượng đây sẽ là một cuộc ẩu đả nảy lửa, một trận chiến gay kịch giữa hai phe để tranh giành trái tim của những cô gái\dots\ thì bạn đã nhầm hoàn toàn. Trận đấu này diễn ra chóng vánh và hoàn toàn một chiều đến mức đáng thương, gói gọn trong vài diễn biến ngắn ngủi sau đây\dots

Zetaro vừa xoa tay vừa hào hứng chuẩn bị tham gia: ``Triển thôi anh em ơi!''

Marine lập tức xoay người, giọng đầy vẻ chỉ huy ra lệnh: ``Để xem bọn chúng có bước qua nổi chúng ta không! Rushia!''

Rushia giật mình vì bị gọi đột xuất, vội vàng đáp lời: ``C-Có!''

``Triệu hồi đám xương của bà lên ngay!''

Không để mất thời gian, nàng Pháp sư gọi hồn lập tức triển khai ma thuật, triệu hồi một dãy bộ xương từ dưới lòng đất ngoi lên nhanh chóng. Những bộ xương xông tới bao vây và ôm ghì chặt lấy đám côn đồ, khiến chúng không thể cử động được dù chỉ một ngón tay.

``\textbf{Cái quái gì thế này!?}'' đám lưu manh gào lên trong cơn hoảng loạn tột độ.

Rushia cũng chẳng ngần ngại, hét lớn ra lệnh tiếp: ``\textbf{Hakimi-san! Noel! Vùng lên đấm chúng nó đi!}''

Noel hừng hực khí thế, nắm chặt bàn tay sẵn sàng chiến đấu: ``Lên liền đây!''

Hakimi bẻ khớp tay phát ra tiếng ``rắc rắc'' đầy ngầu: ``Không cần cô nhắc đâu!''

Lực đấm của Hakimi thực ra cũng không đến nỗi quá khủng khiếp, nhưng tên tay sai bị cậu ấy đấm trúng thì yếu đến mức bất ngờ. Chỉ cần một cú đấm thẳng vào mặt là đủ khiến hắn ngửa bụng, quay lơ nốc-ao ngay tại chỗ.

Còn với nàng hiệp sĩ Noel thì\dots\ nói sao nhỉ? Cú đấm của em ấy lại may mắn không trúng vào người bọn chúng, nhưng chính hành động đó lại khiến cả đám côn đồ sợ đến hồn xiêu phách lạc. Em ấy đã đấm thẳng vào bức tường hành lang, làm vỡ tung cả một mảng gạch xi măng to tướng. Dù là trúng tường, sức mạnh chấn động từ cú đấm đó đủ khiến lũ du côn sợ hãi muốn tiểu tiện ngay tại chỗ.

Chỉ trong nháy mắt, toàn bộ đám tay sai đã bị hạ gục hoàn toàn. Bây giờ chỉ còn lại mỗi tên cầm đầu, vẫn đang cầm lăm lăm cây thước gỗ. Nhìn chứng kiến cảnh tượng kinh hoàng diễn ra quá nhanh, mặt hắn tái nhợt xanh xao, sạch sẽ mọi vẻ hung hăng ban đầu.

Hắn ta lập tức quay ngoắt đầu, vứt phất cây thước gỗ rồi vắt chân lên cổ định chạy mất: ``Cứ chờ đấy! Tao sẽ báo thù!''

Nhưng Flare và Pekora đời nào để hắn thoát dễ dàng như thế. Hắn vẫn đang bị những chiếc chân xương ma thuật của Rushia giữ chặt dưới sàn nhà, không thể nhích được dù chỉ một bước.

Tên côn đồ hốt hoảng gào lên trong tuyệt vọng: ``S-Sao tao không chạy được!?''

Pekora nhếch môi trêu chọc tên đầu trọc đang hoảng sợ: ``Sao lại nhát như thỏ đế thế kia, peko?''

Zetaro đứng bên cạnh vô tình gãi đầu thắc mắc: ``Ơ nhưng mà bạn mới là thỏ mà?''

``Im đi!'' nàng Thỏ lập tức quay sang lườm Zetaro một cái cháy mặt, rồi quay lại dặn Rushia. ``Rushia, nhờ bà đấy. Hôm nay bà là nhân vật chính mà!''

Rushia ngơ ngác không ngờ mình bị đẩy lên trung tâm, vội vàng lắc đầu: ``Đ-Đừng có đùn đẩy tôi chứ-nanodesu! Flare? Pekora?''

Flare và Pekora đồng thanh đáp lời đầy ủng hộ: ``Có bọn mình đây!''

``Cùng nhau hạ gục hắn ta đi!''

Nàng bán yêu tinh Flare lập tức rút ra bộ cung tên theo phong cách Minecraft, chĩa thẳng vào mặt tên cầm đầu. Còn Thỏ xanh Pekora thì nhanh tay lôi ra mấy khối thuốc nổ TNT, cũng từ chính tựa game đó.

Zetaro trợn tròn mắt khi thấy những vật thể kỳ lạ: ``\dots Mấy cái đó bọn nó lấy ở đâu ra vậy?''

``Đừng hỏi tao, tao cũng không biết,'' tôi chỉ biết thở dài đáp lời cậu bạn.

Tên côn đồ run cầm cập, giọng nói lắp bắp vì sợ: ``B-Bọn mày định làm cái gì nữa!?''

Chẳng buồn nghe hắn lảm nhảm thêm câu nào, Flare đã thả tay bắn liền ba mũi tên. Hai mũi đầu trượt qua người, nhưng mũi thứ ba cắm phập thẳng vào bắp chân phải của hắn. Máu bắt đầu rỉ ra, tên đó gào lên đau đớn thảm thiết.

Cùng lúc đó, Pekora đã nhanh chóng xếp các khối TNT ngay dưới chân hắn và cầm sẵn bộ đánh lửa trên tay. Em ấy cúi xuống, giọng đầy vẻ vênh váo hỏi: ``Còn lời gì trăn trối nữa không, peko?''

Tên cầm đầu gào lên đầy giận dữ: ``Tổ sư cái mả cha chúng mày!''

Nàng Thỏ bật cười khoái chí trước lời thề cuối cùng: ``Tao sẽ coi như đó là lời trăn trối cuối cùng của mày!'' Nói rồi, em ấy quẹt lửa thẳng vào khối thuốc nổ.

Pekora lập tức quay lại hét lớn cảnh báo mọi người: ``\textbf{Mọi người, lùi lại phía sau mau! Kuon-senpai, chạy đường khác đi!}''

``R-Rồi!'' tôi giật mình, vội vàng xoay người chạy theo hướng ngược lại.

Tên côn đồ gào lên trong tuyệt vọng khi thấy quả sắp có chuyện: ``Tao nguyền rủa chúng mày, cái đám ngu ngốc này--!''

Chưa kịp để hắn dứt câu, mấy khối TNT nổ tung một tiếng ``ĐÙNG'', hất văng hắn bay thẳng lên không trung. Vụ nổ cũng để lại một khoảng sập lớn, cắt đứt hoàn toàn lối đi ở hành lang. Tôi phải vội vàng chạy ngược xuống tầng dưới, vòng ra phía sau trường để tránh bị kẹt lại trong hiện trường.

Pekora phủi phủi lớp váy dính chút bụi bẩn sau vụ nổ vừa rồi, ngẩng mặt đầy vẻ tự mãn: ``Tưởng mà đến gần bọn tao dễ dàng thế à? Mơ đi cưng!''

Vừa nói, nàng Thỏ vừa vỗ vỗ mấy khối TNT còn sót lại trong tay, như thể vừa hoàn thành một chiến công vĩ đại nhất đời mình. Thật chẳng có gì ngạc nhiên khi em ấy luôn thích thú với mọi thứ nổ tung.

Marine cũng chẳng chịu kém cạnh, chống hai tay vào hông, ngực ưện lên đầy vẻ kiêu hãnh: ``Một khi Thuyền trưởng Houshou Marine đã nhấc chân vào trận, thì đám lưu manh như bọn bây có đáng là cái đầu đinh gì đâu!''

Cậu bạn Zetaro đứng bên cạnh nhảy cẫng lên không trung, hào hứng hùa theo không kém: ``Đúng đó thôi! Đám kia có thấy được sức mạnh khủng khiếp của hội chúng ta chưa hả!?''

Hakimi liền liếc sang cậu bạn thân một cái đầy khinh bỉ: ``Mày có cống hiến được cái méo gì vào trận chiến mà gáy to thế kia? Ở nguyên một đống chỉ biết cổ vũ thì giỏi lắm hả?''

``Đâu có! Tao đóng góp tinh thần mà! Cổ vũ nhiệt tình thế kia còn chưa đủ à!?'' Zetaro vội xua tay bào chữa, mặt méo lại vì bị phốt đúng chỗ yếu.

Flare từ từ quay sang, dội thẳng vào cậu bạn một gáo nước lạnh không thương tiếc: ``Bọn mình đâu có cần cái cổ vũ đó đâu chứ. Mấy việc nhỏ nhặt thế này, bọn mình giải quyết một mình cũng xong xuôi hết.''

Nghe câu nói đó, Zetaro lập tức đứng hình như bị đóng băng, hai tay ôm chặt lấy ngực như thể vừa nhận một phát chí mạng vào tim: ``Hự\dots\ sát thương tâm lý quá nặng nề bạn ơi\dots!''

Nàng Yêu tinh vàng thấy cậu bạn phản ứng mạnh vậy lập tức cuống quýt, vội vàng bước tới vỗ vai an ủi: ``B-Bạn ổn chứ!? Mình xin lỗi nhé\dots\ mình đâu có ngờ bạn sẽ buồn thế này.''

Hakimi chỉ biết thở dài, vỗ vỗ lên vai Zetaro như an ủi một đứa trẻ hư: ``Thôi nào, nó lầy lắm, có sao đâu. Kệ cha nó đi, cứ để nó diễn hết vai diễn của mình đã.''

Rushia bĩu môi nhìn Marine đang hùng hổ vừa nãy, liền nhếch mép chọc ghẹo: ``Mà nói thật thì Marine cũng có làm được gì đáng kể đâu-nanodesu. Cả trận chỉ tát có một cái thì có gọi là chiến công đâu.''

Marine giật nảy mình lên như bị điện giật, vội vàng phản bác ngay không chịu thua: ``Đ-Đâu có! Cái tát đó làm tên kia ngửa cổ ra mà, mạnh lắm chứ bộ!''

``Nếu để tôi giải quyết thì từ đầu tên đó đã nằm dài dưới đất rồi, đâu cần phải phiền phức ra tay làm những trò nhỏ nhặt thế kia,'' nàng Pháp sư thở dài ra vẻ thấm thía, lắc đầu liên tục.

Nàng Thuyền trưởng lập tức hóa đá tại chỗ, mắt tròn thở ra không nói nên lời: ``R-Rushia!? Sao bà có thể nỡ lòng nói mình như thế\dots? Tại sao lại đối xử tàn nhẫn với bạn bè như vậy!''

Flare thấy tình hình bắt đầu căng thẳng lại lập tức hoảng hốt, vội vàng bước sang hòa giải: ``Ơ kìa hai bà này, đừng có cãi nhau nữa! Không lẽ Marine cũng dỗi Rushia thật luôn rồi sao!?''

Hakimi đứng bên cạnh chỉ biết tặc lưỡi liên tục, lắc đầu không tin vào mắt mình: ``Bớt diễn sâu đi mấy bạn ơi, có cần phải kịch tính thế này không hả.''

Khi tôi vất vả vòng ngược lại được vị trí của mọi người sau khi chạy trốn khỏi vụ nổ, thì thấy cả nhóm đang đứng nói chuyện xô xát rôm rả, chẳng có vẻ gì là vừa đánh nhau xong. Zetaro thấy tôi quay liền nghiêng đầu ra hiệu: ``Thôi, dù gì Kuon cũng về tới nơi rồi, ta dời bộ về lớp thôi, kẻo muộn.''

Mọi người lập tức gật đầu đồng loạt, ai nấy đều tán thành: ``Ừ đúng rồi, chứ cứ đứng đây nữa tí nữa bảo vệ trường lên truy bắt thì bọn mình tiêu hết.''

\subsubsection*{\phantomsection\label{ep-C7.4}}
\addcontentsline{toc}{subsubsection}{{\sen{-Partie 4-}} {\caviar Lễ tốt nghiệp của Uruha Rushia}}
\stpt{-Partie 4-}\stcap{Lễ tốt nghiệp của Uruha Rushia}

Và rồi, tiết mục được chờ đợi nhất trong ngày cũng chính thức mở màn: buổi sinh hoạt lớp cuối cùng, một dấu chấm hết cho quãng đời học sinh của chúng tôi.

Dưới ánh đèn vàng ấm áp của lớp học, cô Katsura đã chuẩn bị sẵn những ly trà sữa mát lạnh cho toàn bộ 12A3. Thực ra là gần như cả lớp thôi, vì có mấy tên trốn học về sớm trước đó rồi. Nhưng nhờ vậy mà lại còn dư đúng năm cốc, vừa hay đủ chia cho các thành viên Thế hệ thứ Ba cho buổi học đầu tiên và cuối cùng của họ.

``Còn dư đúng số lượng để chia cho mấy đứa luôn, may ghê,'' cô giáo cười lớn khi đưa từng ly trà đến tay năm cô gái.

Nếu đi theo đúng kịch bản học đường điển hình, buổi học cuối cùng này sẽ là lúc cả lớp ngồi quây quần, rỉ tai nhau những kỷ niệm, những tâm tư chưa kịp nói sau bao năm chung mái trường. Nói về mặt lý thuyết là vậy, chứ tập thể lớp A3 chúng tôi cũng chẳng đoàn kết khăng khít đến mức sướt mướt rơi nước mắt như trong phim.

Dù không có những cảnh tình cảm quá mức, vẫn có những nhóm bạn lần lượt đứng lên chia sẻ cảm nghĩ về những năm tháng cùng học. Có câu chuyện cười đến bể bụng, có kỷ niệm cảm động lắng sâu, và cũng có những câu chuyện khiến tôi phải ngồi lặng người một mình suy ngẫm. Cứ như vậy, cả buổi học tự nhiên biến thành một buổi ``tốt nghiệp mini'' riêng của lớp 12A3.

Trong số tất cả các nhóm, nhóm năm cô gái của nhà Tam Giác Xanh có lẽ là nhóm để lại ấn tượng sâu sắc nhất, khi cả năm người đã thống nhất chọn Rushia làm nhân vật chính trong phần phát biểu của mình. Mỗi tội, bầu không khí buồn thương thì ít mà những pha tấu hài thì vô số, dù phải công nhận các em ấy đã bỏ rất nhiều công sức để chuẩn bị kịch bản chỉn chu từ lâu.

``Và bây giờ, đến lượt nhóm của Thế hệ thứ Ba lên phát biểu cảm nghĩ. Các con lên đi,'' cô Katsura mỉm cười gọi, giọng ấm áp lan tỏa khắp lớp. Ánh mắt cô tràn đầy sự chờ đợi khi nhìn năm cô gái đứng dậy.

``Vâng~!'' Cả năm cô gái đồng thanh reo lên, tiếng nói trong trẻo như chuông gió đập vào nhau.

Họ cùng nhau bước lên bục giảng, bước chân nhẹ nhàng không gây ra tiếng động nào. Pekora tiến ra trước, cất giọng mở lời: ``Mình xin phép có vài điều muốn nói với mọi người, peko.''

Thấy dưới lớp vẫn còn tiếng xì xào bàn tán nhỏ, cô giáo chủ nhiệm đập nhẹ cái thước gỗ xuống bàn theo thói quen: ``Các con trật tự cho các bạn ấy nói nào.'' Kỳ lạ thay, hệt như lúc đầu buổi, cả lớp lập tức im phăng phắc. Được rồi, để xem mấy cô nàng này sẽ chia sẻ điều gì nào. Nhất là khi đây là lần đầu tiên và cũng là lần cuối cùng họ cùng đồng hành với tập thể 12A3 chúng tôi.

``Mặc dù hôm nay mình mới chỉ quen các bạn\dots'' Pekora mỉm cười, gương mặt tròn trịa toát lên sự chân thành. Ánh mắt cô lướt qua từng góc lớp, ghi nhớ mọi khuôn mặt.

Noel tiếp lời, giọng nói trầm ấm như đang kể một câu chuyện cổ: ``\dots nhưng chúng mình cũng đã học được rất nhiều điều cùng với mọi người.'' Cô nàng gật đầu nhẹ, như muốn khẳng định lời mình nói là sự thật.

Flare gật đầu theo, mái tóc vàng óng rủ nhẹ trên vai: ``Chúng mình đã có trải nghiệm tuyệt vời khi được quay trở lại thời trung học, dù thời gian gặp gỡ các bạn có hơi ngẩn ngủi.'' Ánh mắt nàng tràn đầy sự biết ơn.

Marine nháy mắt tinh nghịch, ngón tay chỉ khẽ vào hướng tôi ngồi: ``Mình thấy các bạn cũng khá tốt với bọn mình đấy chứ! Không như ai đó cứ luôn miệng trêu chọc mình\dots'' Cô nàng bật cười nhỏ sau lời nói.

Vừa dứt lời, cả lớp lập tức ngoái đầu quay về phía tôi, bắt đầu xì xào bàn tán. Những tiếng cười khúc khích vang lên nhỏ khắp các dãy bàn. Tôi chỉ biết cười mỉm đáp lại. Kệ họ đi, ``okay I'm fine'', tôi quen rồi.

Rushia nhẹ nhàng nói, giọng nhỏ như làn sương mai buổi sớm: ``Cũng nhờ buổi học hôm nay mà chúng mình biết được trình độ học vấn của bản thân đang ở đâu nữa.'' Em ấy ngại ngùng cúi mặt xuống một chút.

Pekora tổng kết, bước lên một bước để nói to hơn: ``Vậy nên là, mặc dù không có quá nhiều kỷ niệm với các bạn, nhưng dù sao thì\dots''

Cả năm cô gái cùng cúi đầu sâu, đồng thanh cất giọng: ``\textbf{Cảm ơn mọi người rất nhiều!}'' Lời nói vang vọng trong không gian ấm áp của lớp học.

Cả lớp lập tức bùng nổ trong tràng pháo tay giòn giã để đáp lại tấm chân tình của năm cô gái. Tiếng vỗ tay đổ dồn như mưa rơi trên mái nhà. Nhưng phần chính sau đây mới thực sự bắt đầu.

Pekora quay sang nhìn bốn người còn lại, giọng chuyển sang dịu dàng hơn: ``Có rất nhiều điều mình muốn chia sẻ với các thành viên trong thế hệ thứ 3\dots'' Ánh mắt nàng hướng về Rushia với sự trân trọng.

Flare tiếp lời, bước sang cạnh Pekora với nụ cười ấm áp: ``Nhưng, chúng mình muốn dành tình cảm đặc biệt nhất cho Uruha Rushia, một người đồng đội cực kỳ thân thiết của bọn mình.'' Nàng đặt tay nhẹ lên vai Rushia.

Noel mỉm cười dịu dàng, gật đầu đồng tình với hai người bạn: ``Bọn mình đã thảo luận với nhau từ lâu, và quyết định cậu ấy sẽ là nhân vật chính trong buổi lễ tốt nghiệp này.''

Marine hào hứng vỗ tay nhỏ, mắt sáng rực lên như có ngôi sao: ``Mình mong các bạn sẽ chú ý theo dõi nhé. Rushia, bà chuẩn bị chưa?'' Cô nàng sau đó nhún nhảy nhẹ trên chân, trông không thể chờ đợi hơn.

Rushia hơi lúng túng, má ửng hồng vì bị đổ dồn sự chú ý: ``R-Rồi. Kuon-senpai, nhờ anh lấy giúp em mấy cái ghế xếp ở ngoài hành lang vào được không?'' em ấy vội vàng nhờ cậy tôi để che giấu sự ngại ngùng.

``Có liền đây!'' Tôi đáp, nhanh chóng đứng dậy khỏi chỗ ngồi. Tôi cũng muốn giúp đỡ để buổi lễ diễn ra suôn sẻ.

Tôi mờ mờ hiểu được lý do vì sao họ lại chọn Rushia làm ngôi sao chính cho buổi lễ ``tốt nghiệp mini'' này, nhưng tôi quyết định sẽ giữ im lặng. Dù trong thâm tâm, tôi thật sự không muốn nhắc đến những ký ức buồn đó một chút nào.

Tôi lóc cóc bê năm chiếc ghế xếp lên bục giảng, bước chân nhẹ nhàng để không làm phiền ai. Giơ hai ngón tay cái hướng về phía em ấy, làm khẩu hình miệng: ``Cố lên nha!''

Cô Katsura gật đầu mỉm cười, nhìn tôi sắp xếp xong xuôi các chiếc ghế: ``Vậy thì Marine-chan, Pekora-chan, Noel-chan, Flare-chan và Rushia-chan, mời các con. Ai muốn phát biểu trước nào?''

Marine nhẹ nhàng đẩy lưng Pekora để thúc giục nàng bước ra, giọng nhỏ đủ cho mấy cô gái cùng nhóm nghe: ``Pekora-chan đi trước đi.''

Không để ai phải chờ đợi lâu, nàng Thỏ lập tức xung phong giơ tay: ``Ey! Rushia, nhờ bà ngồi cạnh Kuon-senpai nhé?''

``Được rồi\dots'' Rushia ngoan ngoãn gật đầu, bước nhẹ nhàng đến chỗ tôi đã chỉ. Tôi nhanh chóng lùi lại một chút để nhường chỗ cho em ấy, tay khẽ kéo ghế ra cho em ấy ngồi thoải mái. Cả ánh mắt trong lớp đều dồn về phía bục giảng, nơi năm cô gái bắt đầu tiết mục của mình.

Ba người còn lại cũng nhanh chóng ngồi xuống những chiếc ghế xếp tôi vừa mang vào, chỉ còn riêng Pekora đứng dững dưng ở giữa, đối diện trực tiếp với Rushia. Cả lớp đều tò mò chờ đợi nhân vật chính của buổi lễ sẽ có gì chia sẻ, và Rushia đáp lại mọi ánh nhìn bằng một nụ cười ngại ngùng nở rộ trên môi.

Sau khi dọn dẹp giọng cho thật ổn, Pekora mới cất giọng vang lên đủ cho cả lớp nghe: ``Ruu-chan! Bà có nhớ lần đầu tiên chúng ta cùng chơi Minecraft với nhau không?''

Rushia bật cười ngay khi nghe câu hỏi, ký ức cũ ùa về trong đầu: ``Có chứ, tôi nhớ rõ mồn một.''

Hakimi ngồi cạnh tôi vừa xâu chuỗi giấy bút vừa lẩm bẩm như đang tự nói với mình: ``\dots Ok, tôi vẫn chưa hiểu chuyện gì đang xảy ra lắm.''

Tôi khẽ huých vai cậu bạn, rồi thầm thì nhỏ đủ chỉ có hai người nghe: ``Người ngoài cuộc thì không hiểu được đâu. Cứ im lặng theo dõi đi.''

Pekora chẳng để ý đến những tiếng xì xào nhỏ dưới lớp, tiếp tục câu chuyện của mình với giọng càng lúc càng ấm áp: ``Hồi đó dù chúng mình chưa thực sự hiểu ý nhau lắm, lúc nào cũng có người chết lạc hướng giữa server, nhưng khoảng thời gian đó vui cực kỳ luôn đó!''

``Ừ, tôi nhớ mà.'' Rushia gật đầu liên tục, mắt cũng ánh lên niềm vui khi nhớ lại những kỷ niệm cũ.

``Nhưng không lâu sau, chúng ta đã trở thành những người đồng nghiệp thấu hiểu nhau, cùng nhau xây dựng nên bao nhiêu công trình to lớn trên thế giới ảo, và rồi tình cảm giữa chúng mình ngoài đời cũng ngày càng thắm thiết hơn!''

``Đúng thế thật!'' Rushia đáp lời, giọng cũng nghẹn ngào không kém bạn mình.

Pekora rụt rè lau khóe mắt, giọng run run như sắp khóc: ``Tôi nên nói thế nào nhỉ\dots\ Cảm ơn bà đã luôn đồng hành cùng tôi! Cảm ơn bà rất nhiều!''

Nàng Pháp sư đáp lại bằng một nụ cười vô cùng rạng rỡ, đủ để xóa tan mọi nỗi buồn. Cả lớp dù chẳng hiểu câu chuyện Minecraft hay server kia là gì nhưng cũng vỗ tay rầm rầm theo thói quen, không khí trong lớp ấm áp lạ thường.

Tuy nhiên, nàng Thỏ chẳng muốn bầu không khí trở nên quá sến sẩm và ngột ngạt, nên nàng cố gắng nghĩ ra điều gì đó tấu hài để phá vỡ: ``Mà\dots\ cái này nói sao cho hết ý nhỉ?''

Boji, người ngồi ở dãy bàn đầu, liền ngẩng mặt lên thắc mắc: ``Ờ thì\dots\ giống như những gì bạn vừa nói thôi mà?''

Kihataru ngồi cạnh cậu ấy liền phán một câu xanh rờn: ``Thiếu mỗi từ `peko' ở cuối nữa là chuẩn bài rồi.''

``K-Không phải thế mà, peko!'' Pekora đỏ bừng mặt chống chế, tay vẫy vẫy giải thích. ``Có ai biết tôi nên nói thế nào cho chuẩn không?''

Marine đứng bên cạnh liền thì thầm nhắc bài, giọng nhỏ như muỗi kêu: ``Pekora, Pekora! Đây này: `Cảm ơn bà vì mọi thứ!' ''

``À à phải rồi! Cảm ơn bà vì đã đồng hành cùng tôi!'' nàng Thỏ lặp lại y nguyên câu đó, chẳng thay đổi được bao nhiêu.

Hakimi ngồi dưới lớp tặc lưỡi khó hiểu: ``Có khác méo gì câu trước đâu\dots''

Bất chấp mọi lời phàn nàn, Pekora vẫn hô lớn để đếm nhịp cho cả nhóm: ``Được rồi, hai, ba\dots''

Cả năm người họ cùng reo lên: ``\textbf{Rushia-chan~!}'' một cách vô cùng lệch nhịp và hỗn loạn, nghe như một đám học sinh đang cổ vũ trong buổi thể thao. Rushia bật cười ngặt nghẽo trước hành động ngốc nghếch của mấy người bạn: ``Cái quái gì thế này!?''

Kijiru ngồi cùng dãy với tôi cũng ngơ ngác lắc đầu: ``Tôi cũng đang định hỏi câu đó đấy\dots''

Tôi đành phải đứng lên giải thích cho cả lớp hiểu về đặc trưng kỳ lạ của nhóm này: ``Đặc trưng của nhóm bọn họ là vậy đấy mọi người. Này mấy đứa, làm thế nào để mọi người cùng cười vui vẻ được không?''

Nàng Thỏ đỏ bừng mặt vì ngượng, vội vã ngồi thụp xuống ghế trước tiếng cười tủm tỉm của cả lớp: ``Nhưng em không biết làm cách nào hết, peko! À thì\dots\ người tiếp theo đi!''

Đến lượt Flare bước lên trung tâm bục giảng, nối tiếp câu chuyện của Pekora. Suốt buổi lễ này, cứ mỗi lần có người nói gì đó, Rushia đều đáp lại bằng những tiếng ``ừ'' hay ``rồi'' nhỏ nhẹ như một phản xạ tự nhiên, một thói quen khó bỏ từ bao giờ.

Cô Katsura mỉm cười khuyến khích, giọng nói ôn tồn như đang dỗ dành một đứa con gái nhỏ: ``Shiranui-san, con có điều gì muốn chia sẻ với bạn ấy không?''

``Có ạ,'' nàng Yêu tinh vàng khẽ cúi đầu lễ phép, sau đó ngước mặt nhìn quanh lớp với ánh mắt hơi bỡ ngỡ. ``Xin phép mọi người\dots''

Koheki không kiềm được mà cất tiếng giục, giọng cười khúc khích làm bầu không khí bớt ngượng ngùng: ``Thì bạn cứ nói đi, có ai thúc ép đâu mà rụt rè thế.''

Tôi liền huých nhẹ vai cậu bạn Hóa học, nhắc nhỏ nhẹ rồi hướng tới bục giảng: ``Đừng làm người ta căng thẳng chứ\dots\ Em cứ nói tự nhiên như lúc mấy đứa hay nói chuyện với nhau là được.''

Flare hít một hơi thật sâu để lấy lại bình tĩnh, sau đó mới quay mặt nhìn thẳng về phía Rushia. ``Được rồi\dots\ Rushia, cái hồi lần đầu tiên tôi với bà làm buổi stream uống rượu buổi tối cùng nhóm NoeRushiFure ấy\dots''

Nghe đến hai chữ ``uống rượu'', cả lớp lập tức ngơ ngác trố mắt, không ai tin được vào tai mình. Đám bạn bắt đầu xì xào bàn tán với đủ thứ suy đoán trái chiều: ``Trông ngoan hiền lễ phép thế mà cũng là tay uống rượu gớm nhỉ!''

Nhưng xin thưa, Flare vẫn chưa là gì so với nàng Hiệp sĩ tâm hồn cơm bò Noel kia đâu. Mấy cô nàng này, chuyện uống rượu nói chuyện đêm khuya là thói quen thường nhật rồi. Hakimi há hốc mồm đến mức suýt rơi cả cặp kính trên mũi, không giấu được sự bất ngờ: ``Vãi chưởng thật\dots''

Zetaro cũng lắc đầu ngao ngán, vỗ vai cậu bạn bên cạnh như vừa chứng kiến một điều gì đó phi thường: ``Mấy em này uống được cả rượu cơ à\dots\ Bái phục, còn khỏe hơn cả anh em mình.''

Tôi liền tạt một gáo nước lạnh vào đám người đang phóng đại vấn đề, cười khẩy để bình ổn tình hình: ``Chắc gì người ta đã uống rượu thật.''

Flare vội vàng ngượng ngùng xua tay đính chính, má cô nàng ửng hồng vì hiểu lầm của mọi người: ``Đ-Đúng đó! Tôi đang nói đến đâu rồi nhỉ\dots\ *e hèm* Tôi cứ tưởng lúc đó bà mới chỉ 16 tuổi thôi, thế mà sau này tôi mới biết là bà tận 1600 tuổi lận đấy! Rushia-chan!''

Ngoại trừ Rushia chỉ biết cười trừ trước câu nói dí dỏm của cô bạn, cả lớp lại thêm một trận há hốc mồm kinh ngạc. Tiếng than thở vang lên khắp các dãy bàn, ai cũng không thể tin vào con số vừa nghe: ``Bạn ấy lớn tuổi đến thế cơ á!?''

Tất nhiên, nàng Pháp sư ngượng đến mức mắng yêu cô bạn cùng thế hệ, tay khẽ đập nhẹ vào ghế ngồi: ``Đừng có nói kiểu như tôi gian lận tuổi tác thế chứ!''

Khan lắc đầu tặc lưỡi, vẫn chưa hết bất ngờ trước thân thế bí ẩn của Rushia: ``Bây giờ tôi mới biết đấy\dots''

Rushia lấy hai tay che lấy mặt, rên rỉ vì xấu hổ khi cả lớp đều đang đổ dồn ánh nhìn về phía mình: ``Bây giờ mọi người biết hết rồi sao\dots\ Cái gì vậy chứ!''

Hakimi không kiềm được mà thốt lên câu tục chửi quen thuộc, cố gắng che giấu sự sốc của mình: ``Vãi cả beep\dots'' Tôi nhếch môi cười nhỏ, quay sang nói với Hakimi để cậu bạn đừng quá ngạc nhiên: ``Thế là vẫn còn bình thường chán. Tôi còn gặp mấy ca có số tuổi oái oăm hơn nhiều.''

Cậu bạn liền tò mò xoay người sang, ánh mắt hiện rõ sự hiếu kỳ: ``Kiểu gì?''

``Ông có tin công ty tôi làm có thành viên 0 tuổi mà đã biết lập trình máy tính siêu việt không?''

Hakimi trợn tròn mắt, mồm há rộng không thể khép lại được: ``Vãi sít\dots''

``Cứ bình thường đi anh bạn.'' Tôi quay sang Flare, ra hiệu cho cô nàng tiếp tục chương trình. ``Flare-chan, xong chưa em?''

``Rồi ạ!'' Flare mỉm cười vui vẻ, sau đó lùi lại để nhường chỗ cho người tiếp theo. ``Noel, tới bà đấy!''

Cô Katsura gật đầu lơ đãng, giọng còn vương chút ngỡ ngàng sau màn thông tin vừa rồi: ``Ờm\dots\ vậy đến lượt người tiếp theo lên đi.''

Noel đứng dậy, vẫy tay chào cả lớp với nụ cười hiền hậu: ``Ừm, đến tôi đây! Rushia, cậu còn nhớ buổi collab đầu tiên chúng ta cùng phát sóng không? Hồi đó cả hai đứa đều ngơ ngơ ngáo ngáo, làm gì cũng bỡ ngỡ lắm\dots''

Nàng Pháp sư cười khúc khích, gật đầu đồng tình: ``P-Phải ha, hồi đó đúng là dở hơi thật sự\dots''

Nàng Hiệp sĩ tiếp tục câu chuyện với giọng ấm áp, đôi mắt nhìn xa xăm như đang nhớ lại kỷ niệm cũ: ``Lúc đầu chúng mình cũng chẳng biết nên bắt đầu chuyện gì, cứ im lặng ngượng ngùng suốt buổi, cơ mà\dots''

Rushia lập tức cắt ngang, giọng có chút phì cười trước lời bạn: ``Bà thật sự nghĩ thế sao!?''

Zetaro ngồi cuối lớp nghiêng đầu thì thầm với người bên cạnh, giọng nhỏ như muỗi kêu: ``\hl{Nghe mô típ này sao quen thế, giống hệt ai đó mới đến lớp hồi đầu năm ấy nhỉ?}''

Koheki ngồi cạnh liền lườm một phát, tay vừa đạp nhẹ chân cậu bạn: ``Im mồm đi, nói nhiều quá.''

Cả lớp bật cười rộn rã trước cảnh hai cậu bạn cãi nhau nhỏ. Nàng Hiệp sĩ tóc bạc lại tiếp tục câu chuyện của mình, giọng dịu dàng làm dịu bầu không khí: ``Mình cứ tưởng chúng ta sẽ giữ không khí ngượng ngùng ấy mãi cơ đấy!''

Rushia gật đầu lia lịa, nhớ lại cảnh tượng hồi đó vẫn còn rõ mồn một: ``Ừ, đúng đúng, hồi đó tôi cũng lo lắm.''

``Nhưng mà, ngay đêm hôm ấy, Rushia đã nhắn tin riêng hỏi mình rằng `liệu những gì tôi nói trên stream có ổn không nhỉ?' các kiểu.''

Zetaro bất ngờ giơ tay ngắt lời, mặt tò mò không giấu được: ``Xin lỗi hỏi không liên quan, nhưng bạn ấy hỏi về cái gì thế?''

Noel bỗng dưng ấp úng, má ửng hồng vì ngượng, vội vàng tìm cách diễn đạt khác: ``À, cậu ấy muốn biết làm cách nào để cho\dots\ cái núi ấy nhô cao lên\dots\ à không, ý mình là là bóp\dots''

Cậu bạn tròn xoe mắt, càng nghe càng mơ hồ: ``Bóp gì cơ?''

``Cái\dots\ cái núi ấy! Chính cái núi thôi!'' Noel càng nói càng vội, mặt đỏ bừng lên.

Zetaro vẫn không hiểu, lặp lại câu hỏi với giọng ngây thơ: ``Núi gì nữa vậy?''

Nàng Hiệpsix cuống cuồng xoay tay, cố gắng giải thích mà không làm lộ chuyện tế nhị: ``Núi\dots\ ơ\dots\ Bạn cũng hiểu mình đang nói đến cái gì đúng không!?''

Zetaro lắc đầu nguầy nguậy, thật sự không mường tượng ra được: ``Nhưng mà tôi không hiểu thật sự!''

Noel rối bời đến mức kéo cả mái tóc của mình, giọng lạc đi vì ngượng: ``Nói thế nào cho bạn hiểu nhỉ\dots\ Ờm\dots''

Rushia bắt đầu nổi đóa, mặt hầm hầm quát lên khi thấy Noel càng nói càng rối: ``Này, bà đang lắp bắp cái gì thế hả!? Cứ nói thẳng ra đi!''

Tôi vội vàng can ngăn, vỗ vai cậu bạn lắm lời để giải vây cho hai cô gái: ``Zetaro-senpai, đừng có trêu em ấy nữa, coi chừng bị ăn đòn bây giờ.''

Hakimi cũng phụ họa, tay chỉ thẳng vào miệng cậu bạn ra hiệu ngậm lại: ``Mày ngậm cái mồm lại cho đời nó bình yên đi, không là bọn tao không giúp được mày đâu.''

Zetaro vẫn ngoan cố, quay sang hỏi cả dãy bàn phía sau để tìm đồng minh: ``Nhưng thật sự anh em có ai hiểu bạn ấy đang nói gì không? Chứ tôi thật sự không hiểu!''

Tới lúc này, cậu bạn thở dài bất lực, quay sang gọi cậu bạn chuyên Hóa bên cạnh: ``\dots Koheki ơi, tẩn cho nó một trận đi, để nó hết cái tật lắm lời này.''

Chưa để cậu bạn Bách khoa kịp phản ứng, Koheki đã lao tới ``úp sọt'' ngay, đè đầu cậu bạn xuống sàn lớp để dạy dỗ. Đúng lúc đó, Marine đứng bên cạnh bỗng phang một câu thẳng thắn không ngại ngần: ``Chắc là Noel muốn làm cho bộ ngực của mình đầy đặn hơn chứ gì nữa! Cái ``núi'' đó mà!''

Cả lớp lập tức bật cười nghiêng ngả, tiếng cười vang vọng khắp phòng học. Zetaro đang bị đè dưới sàn cũng reo lên như vừa khai phá ra bí ẩn: ``À~ Ra là thế! Tôi hiểu rồi!''

Rushia mặt lập tức biến sắc, sát khí toả ra quay sang mắng Marine: ``\dots Bà muốn chết đúng không, hả Marine!?''

Nàng Thuyền trưởng hoảng hốt xua tay như muối bỏ bể, vội vàng thanh minh: ``K-Không phải thế! Tôi chỉ đang giải thích giùm Noel với cậu bạn kia thôi mà! Sao tôi lại bị vạ lây thế này!''

Cậu bạn Hóa ho vừa giữ chặt Zetaro không cho cậu vùng dậy, vừa hỏi với giọng nghi ngờ: ``Mày giả vờ không hiểu hay cố tình không hiểu đấy hả Zetaro? Cái chuyện rõ ràng thế mà cũng hỏi.''

Zetaro mếu máo dưới lớp sàn, giọng oan ức lắm: ``Tao không hiểu thật mà! Ai biết các cô ấy nói ẩn dụ thế đâu!''

``Ấn nó xuống thêm tí nữa đi Koheki, để nó hết cái tật ngốc nghếch này,'' tôi bồi thêm một câu, cười khúc khích nhìn cảnh tượng.

``Ơ từ từ-- \textbf{Á á á!! Dừng tay đi mấy thằng bạn!}'' tiếng Zetaro thất thanh vang lên khi bị đè mạnh hơn.

Khan ngồi phía trước lắc đầu chán nản, quay lại lườm cậu bạn một cái: ``Bất lịch sự vãi chưởng\dots\ Ông cứ im mồm đi cho bọn tôi nhờ cái, đừng làm rối tung buổi lễ lên.''

Rushia vẫn còn hừm mũi, lườm Zetaro với ánh mắt không mấy thiện cảm: ``\dots Tí nữa tôi dùng ma thuật kéo cái đứa đang bị úp sọt kia xuống thẳng địa ngục luôn, cho biết mùi.''

Tôi vội vàng vỗ vai em ấy để làm hòa, giọng nhẹ nhàng bình ổn: ``Bình tĩnh nào Ru-chan\dots\ Marine-chan với Zetaro cũng chỉ đùa thôi, không có ý xấu đâu mà. Phải không Noel-chan?''

Noel vội vàng chốt lại vấn đề trước khi tình hình thêm rối, gật đầu lia lịa: ``D-Dù sao thì! Mình chỉ muốn nói rằng mình rất vui vì chúng ta ngày càng gắn kết hơn thế này! Cảm ơn cậu rất nhiều, Rushia-chan!''

Cả lớp liền dành cho nàng Hiệp sĩ một tràng pháo tay giòn giã, tiếng vỗ tay vang lên đầy yêu thương. Chỉ có một cậu bạn đang bị đè dưới sàn còn ấp úng van xin: ``Mau\dots\ mau cho tao đứng dậy đi, tao chịu hết nổi rồi!''

Marine cười trừ nhìn cảnh lộn xộn, nghiêng đầu nói với mọi người: ``Mọi người bây giờ có vẻ hơi khó ăn nhập với nhau nhỉ? Bình thường cũng có hỗn loạn thế này đâu.''

Rushia thở dài xoa thái dương, mặt vừa bực mình vừa bất lực: ``Tôi cũng chẳng hiểu chuyện quái gì đang xảy ra nữa. Nhưng mà\dots\ đúng là cảm giác đây thật sự giống một buổi lễ tốt nghiệp của mình ấy\dots''

Tôi gật đầu đồng tình, giọng trầm ấm nhìn khắp lớp học thân thương: ``Ờ thì, dù sao đây cũng là buổi học cuối cùng của chúng ta rồi mà. Những kỷ niệm lộn xộn thế này mới nhớ lâu chứ.''

Boji ngồi ở dãy bàn đầu tiên mỉm cười dịu dàng, ánh mắt ấm áp dạo qua khắp lớp như đang ngắm nhìn một khung cảnh thân thương. Gương mặt cậu ấy toát lên sự yên bình hiếm có giữa buổi tiễn biệt sắp tới. ``Không sao đâu, mình thấy mọi người cứ vui vẻ thoải mái với nhau như thế này là tốt rồi. Đó mới là điều quan trọng nhất.''

Khan ngồi cạnh cậu bạn chợt nhắc nhở, giọng vừa đủ cho nhóm bạn nghe thấy, ngón tay khẽ chỉ về phía bục giảng. Cậu ấy dường như nhận ra điều gì đó mà phần lớn mọi người chưa để ý. ``Hình như Marine-san vẫn còn điều gì muốn tâm sự với bạn ấy đúng không?''

Vừa dứt câu, Thuyền trưởng lập tức như được nhấc bổng bởi một luồng năng lượng vô hình, nhảy tót đứng hẳn lên trên chiếc ghế xếp cao nhất. Hai chân cpp nàng vững vàng như đang đứng trên boong tàu của mình, hít một hơi thật sâu để lấp đầy cả lá phổi. Rồi dõng dạc gào toàn bộ lực về phía Cô đồng Uruha Rushia đang cười lăn cười bò trên ghế: ``\textbf{Rushia-chan à!!}''

Cú hét bất ngờ vang lên như một tiếng pháo hoa giữa đêm yên tĩnh, làm Rushia giật nảy mình. Em ấy suýt nữa thì trượt khỏi ghế, tay vội vàng bám chặt vào mép ghế để giữ thăng bằng. Ánh mắt đỏ tròn xoe đầy ngỡ ngàng nhìn về phía người bạn thân thiết.

Cô Katsura ngồi ở bàn giáo viên cũng giật thót tim hoảng hốt, tay vội vàng đặt lên ngực như thể vừa bị nhát chết. Cô đứng dậy định chạy tới kéo Marine xuống, giọng nói lắp bắp vì lo lắng. ``X-Xuống đi con ơi! Đứng thế ngã bây giờ!''

Tôi vội vàng giơ tay ngăn cản cô giáo, cố gắng giữ bình tĩnh trước cảnh tượng hài hước đang diễn ra. Tôi hiểu rõ tính cách của các cô gái ấy hơn ai hết, nên biết rằng đây chỉ là một phần của màn trình diễn. ``Kệ bạn ấy đi cô ạ, cứ để cho bạn ấy nói đi. Hẳn là bạn ấy đang có điều gì muốn giãi bộc từ tận đáy lòng đấy ạ.''

Kijiru ngồi dưới bục giảng không chịu nổi nữa, lập tức hò reo cổ vũ như đang ở một buổi hòa nhạc. Cậu ấy vỗ tay rầm rầm, mặt rạng rỡ hẳn lên vì hào hứng. ``Thể hiện đi bạn ơi!''

Nàng tóc nâu đỏ hít một hơi thật sâu, ngực nâng lên như một chiếc buồm căng gió trước khi cất giọng tiếp tục. Các thành viên còn lại của Thế hệ thứ Ba cũng chủ động lui xuống đứng ở hai bên bục giảng, tạo thành một khung cảnh độc đáo. Tất cả ánh mắt trong lớp đều đổ dồn về phía Marine.

``\textbf{Bà là mối tình đầu của tôi đó!}''

Cả lớp lập tức im phăng phắc, không một tiếng động nào có thể được nghe thấy trong không gian lớp học 12A3. Trên đầu mấy đứa bạn tôi mọc ra một đống dấu chấm hỏi to tướng, biểu cảm ai cũng như bị đóng băng trong trạng thái ngỡ ngàng. Đáng chú ý nhất là biểu cảm của hai tên cùng nhóm - Hakimi và Khan, ``top 1 server'' của lớp.

Hakimi đổ mồ hôi hột, trán lấm tấm những giọt mồ hôi lạnh vì sốc quá mức. Cậu ấy quay sang nhìn tôi với ánh mắt không thể tin được, mồm há rộng không thể khép lại: ``Holy shit, phát cơm chó đấy à?''

Khan cũng hoảng không kém, mặt cắt không còn giọt máu khi nghe lời thú nhận bất ngờ. Cậu ấy lắc đầu liên tục như thể muốn xác nhận rằng mình không phải đang nghe lầm: ``Lại còn là yuri nữa chứ!''

Kijiru gãi đầu ngơ ngác, gương mặt đầy thắc mắc sau khi vừa cổ vũ nóng bỏng. Cậu ấy chưa bao giờ tưởng tượng được tình tiết lại đi đến mức này, hoàn toàn nằm ngoài dự đoán: ``Không ngờ tới luôn đấy\dots''

Tôi nhếch môi cười khẩy, hoàn toàn không bị cảnh tượng đó làm cho sốc vì đã quá quen thuộc. Những màn tấu hài kiểu này là một phần không thể thiếu trong cuộc sống hàng ngày của công ty tôi. ``Tôi đã bảo rồi, mấy cái này ở công ty tôi là chuyện như cơm bữa\dots\ Ờ, Marine-chan, tiếp tục đi em. Có gì cứ bộc lộ hết ra.''

Marine hào hứng tiếp lời, không để ý đến những tiếng bàn tán nhỏ bắt đầu nổi lên dưới lớp. Nàng tiếp tục màn tấu hài của mình với thái độ hết sức nghiêm túc, làm cho tình hình càng thêm vui nhộn: ``\textbf{Kể từ lần đầu tiên gặp bà, tôi đã luôn nghĩ rằng bà là người con gái dễ thương nhất quả đất!}''

Rushia gào lên đầy mãnh liệt, tay giận dỗi đập nhẹ vào ghế ngồi. Em ấy không thể tin được vào những lời nói dối trắng trợn của người bạn, mặt đỏ bừng lên vì ngượng: ``\textbf{Nói xạo!!!}''

Hakimi quay sang thì thầm vào tai tôi, giọng đầy lo lắng như đang chứng kiến một thảm họa sắp xảy ra. Cậu ấy chưa từng thấy một màn trình diễn nào kỳ lạ như thế này trong đời: ``\hl{Này Kuon, liệu kịch bản có ổn thật không đấy?}''

Khan tặc lưỡi, đồng ý với những lo lắng của Hakimi với một cái gật đầu nhẹ. Cả hai người đều cảm thấy có gì đó không ổn với chuỗi sự kiện đang diễn ra: ``Nghe nó cứ fake fake thế nào ấy\dots''

Koheki phán một câu xanh rờn từ phía sau, phá vỡ sự im lặng ngượng ngùng giữa tôi và hai người bạn. Cậu ấy không ngại ngần bóc phốt sự gượng gạo của màn kịch: ``Gượng vãi l\dots''

Mặc kệ những lời bàn tán xô xát dưới lớp, cặp đôi hoàn hảo trên bục giảng vẫn tiếp tục màn đối đáp tấu hài của mình. Họ như hai diễn viên chuyên nghiệp đã cùng diễn xuất hàng trăm lần, ăn ý đến mức lạ thường. Chưa kịp để Senchou nói hết câu, ba cô gái còn lại đứng hai bên đã đồng thanh hô lớn vì tưởng em ấy kết bài: ``\textbf{Rushia-chan~!}''

Hô xong mới nhận ra mình bị hớ, cả ba cô gái đều đứng im người ngỡ ngàng vì đã phá hỏng dòng chảy của câu chuyện. Noel là người đầu tiên phản ứng, gương mặt đầy bối rối không biết phải làm gì tiếp theo. ``À, chưa tới lượt hả?'' nàng Hiệp sĩ PON ngớ người.

Chính Thuyền trưởng cũng phải xua tay ra lệnh cho mọi người đứng lùi xuống, không để tình hình thêm rối ren. Cô vừa cười vừa lắc đầu, không thể ngờ được các bạn mình lại hớ đến thế. ``Tôi vẫn còn vài điều muốn nói nữa mà\dots\ Còn nữa mà!'' rồi bật cười hỉ hả. Ba cô gái ngượng ngùng bước lại vị trí cũ.

Cô Katsura dở khóc dở cười nhìn cảnh tượng lộn xộn trên bục giảng. Cô đã dạy học nhiều năm nhưng chưa bao giờ chứng kiến một buổi chia tay cuối lớp kỳ lạ đến thế. ``Ờm\dots\ tiếp đi con.''

``Dạ, thưa cô,'' Marine quay sang Rushia, hạ giọng ấm áp sau khi đã ổn định tình hình. Nàng quay trở lại với phần phát biểu của mình, cố gắng mang lại chút tình cảm chân thành giữa những pha tấu hài. ``Giọng của bà rất trầm ấm, bà lúc nào cũng quan tâm tới tôi\dots''

Rushia mỉm cười ngại ngùng, má ửng hồng vì những lời khen của người bạn. Em ấy cúi mặt xuống một chút, không dám nhìn thẳng vào mắt Marine trước sự chú ý của cả lớp. ``Thôi mà\dots''

``\textbf{Bà cũng cực kỳ dễ thương nữa!}''

Kijiru bóc phốt ngay lập tức, không để cho lời nói dối đó trôi qua một cách dễ dàng. Cậu ấy nhận ra rằng Marine đã lặp lại chính câu nói đó cách đây không lâu: ``Cái này bạn vừa nói ban nãy rồi mà\dots''

Nàng Thuyền trưởng tiếp tục phang tới bến, không hề bị lời bóc phốt của Kijiru làm cho nản lòng. Nàng càng tấu hài càng hưng phấn, quyết định phải đưa màn kịch của mình đến cao trào: ``\textbf{Vậy nên cho dù bà không có ngực đi chăng nữa\dots}''

Rushia lập tức nổi đóa quát lên, mặt hầm hầm như sắp bốc cháy vì lời nói đùa quá trớn của Marine. Em ấy đứng dậy định lao tới đánh người bạn, tay nắm chặt thành quyền: ``\textbf{Im đi!}''

Cả tôi lẫn Hakimi cùng đồng thanh lên tiếng, không thể ngồi yên nhìn tình hình càng thêm điên rồ. Chúng tôi vừa cười vừa lắc đầu, không thể ngờ được Marine lại dám xoáy sâu vào điểm yếu của Rushia: ``Có nhất thiết phải xoáy sâu vào nỗi đau của người ta thế không vậy\dots''

Marine khoanh tay tuyên bố dõng dạc, kết thúc màn tấu hài của mình với một câu nói mà nàng cho là vô cùng sâu sắc. Nàng đứng vững vàng trên ghế, mặt đầy vẻ nghiêm túc như đang thực hiện một lời thề. ``\textbf{Tôi vẫn yêu bà nhiều lắm!}''

Cả lớp lại được trận cười nghiêng ngả trước lời thú nhận vô cùng thẳng thắn của Thuyền trưởng, tiếng cười vang vọng khắp phòng học như một bản nhạc vui. Trong khi đó, Rushia thì phồng má giận dỗi, không nói nên lời trước sự tinh nghịch của người bạn.

``Bạn ấy có khác quái gì tao đâu,'' tên Zetaro lẩm bẩm từ cuối lớp, cuối cùng cũng tìm được điểm chung với Marine. Cậu ấy vừa nói vừa cười, nghĩ rằng mình đã tìm được người đồng điệu. Chưa cần tôi phải cất lời, Koheki đã tiện tay ấn đầu cậu ta xuống bàn.

Rushia bật cười trừ, cơn giận của em ấy cũng tan biến theo tiếng cười của cả lớp. Em ấy lắc đầu bất lực, không thể nào giận được người bạn thân thiết của mình. ``Cái gì vậy trời!?''

Marine nắm lấy tay bạn, bỏ đi hết những giọt đùa nghịch để thay vào đó là một ánh mắt chân thành. Cô nhìn thẳng vào mắt đối phương, giọng nói trở nên ấm áp và đầy tin tưởng. ``Tôi muốn chúng ta hãy mãi mãi ở bên nhau nhé!''

Nàng Cô đồng bất chợt khựng lại, nụ cười trên môi dịu xuống khi cảm nhận được sự chân thành trong lời nói của Marine. Em ấy ngước đôi mắt long lanh về phía người bạn, như thể không thể tin được vào những gì mình vừa nghe. ``Thật ư\dots?''

Tôi không ý kiến gì vì đây là sân khấu của họ, tôi chỉ là một người quan sát bình thường trong buổi tối quan trọng này. Nhưng vài đứa trong lớp đã bắt đầu ngơ ngác trước sự thay đổi cảm xúc xoay như chong chóng của Rushia, không theo kịp được dòng chảy cảm xúc. ``Sao lại chuyển cảnh nhanh thế được\dots'' một ai đó lẩm bẩm nhỏ.

Marine tiếp tục đứng trên ghế gào to, quên mất rằng mình không cần phải hét lên để cả lớp nghe thấy. Em ấy vẫn còn nhiều điều muốn nói, muốn chia sẻ với tất cả mọi người trong lớp. ``\textbf{Dạo gần đây ấy--}''

Tôi vội xua tay can, ngăn cản cô ấy tiếp tục hét lên như thể đang ở trên một con tàu ở giữa biển khơi. Tôi không muốn cả lớp bị điếc tai vì màn phát biểu này. ``Em không cần phải hét toàn bộ lực thế đâu Marine-chan. Bọn anh chưa điếc đâu mà.''

Kihataru gật đầu phụ họa, hoàn toàn đồng ý với những gì tôi vừa nói. Cậu ấy cũng cảm thấy rằng Marine đang hơi quá đà với cách diễn đạt của mình: ``Thật đấy. Mọi người lại tưởng bạn đang tấu hài mất. Với cả xuống ghế đi, đứng bình thường nói chuyện là được rồi.''

Nghe những lời bình luận từ đám ``Rách'' chúng tôi, Thuyền trưởng ngượng ngùng gãi đầu, nhận ra rằng mình đã làm hơi quá. Cô nhanh chóng leo xuống ghế để đứng cho đúng mực, mặt đỏ bừng lên vì xấu hổ. ``Đ-Đúng ha!'' Em ấy hít một hơi sâu rồi tiếp tục.

``Tôi đã suy nghĩ rất lâu về cách để chúng ta có thể luôn ở bên nhau. Vậy nên, Rushia-chan à\dots''

Như đã hẹn trước từ lâu, Flare, Pekora và Noel bước lên bục giảng với sự đồng bộ hoàn hảo. Ba người bạn đã chuẩn bị sẵn sàng cho thời điểm này, cùng hô to với hết sức mình: ``\textbf{Rushia-chan~!}''

Tưởng như bài phát biểu của Thế hệ thứ Ba đã kết thúc ở đó, mọi người đã chuẩn bị vỗ tay cho màn trình diễn, nhưng Marine bất ngờ thêm vào một cái kết hoàn toàn ngoài dự đoán. Nàng vừa cười vừa chỉ về phía tôi, làm cho mọi sự chuẩn bị của các cô gái trở nên vô nghĩa: ``\dots và cả Kuon-senpai nữa\dots''

Tôi giật mình chỉ vào mặt mình, không thể tin được rằng mình lại bị lôi vô tình vào câu chuyện. Tôi đứng ngẩn người, không hiểu tại sao mình lại có mặt trong màn kịch này: ``Ủa, có cả anh nữa à?''

Flare mỉm cười nhìn tôi, ánh mắt đầy ấm áp và thân thiện như thể tôi thực sự là một phần của nhóm. Cô nói với tất cả mọi người trong lớp, giọng nói chân thành lan tỏa khắp phòng. ``Chứ còn sao nữa ạ. Anh cũng là một phần trong tập thể này mà.''

Pekora giơ tay, tiếp lời Flare với nụ cười rạng rỡ thường ngày của nàng Thỏ. Mọi cảm giác ngượng ngùng ban đầu đã tan biến, thay vào đó là sự tự tin và gắn kết. ``Dù chúng ta có đi đâu\dots''

Noel tiếp lời, hoàn thành câu nói mà cả nhóm đã cùng nhau chuẩn bị từ lâu. Nàng nhìn về phía Rushia với ánh mắt đầy yêu thương và trân trọng. ``\dots hay bất cứ nơi nào\dots''

Marine hô lớn, kết thúc câu chuyện với một thông điệp mà cả nhóm đều muốn gửi gắm. Cô nàng muốn mọi người trong lớp đều hiểu được tình cảm đặc biệt của nhóm Thế hệ thứ Ba: ``\textbf{Chúng ta vẫn sẽ mãi là những người bạn vĩnh cửu! Mọi người thấy thế nào!?}''

Cả nhóm đồng thanh hưởng ứng, vung cao tay lên trời như những đứa trẻ đang chơi đùa trong công viên. Họ hoàn toàn thoải mái thể hiện tình cảm của mình trước toàn bộ lớp học. ``\textbf{Đồng ý!}''

Tôi vẫn đứng ngơ ngơ ngác ngác giữa bục giảng, không hiểu chuyện gì đang xảy ra với mình. Tôi bị lôi vô thức vào một màn kịch mà mình không hề được chuẩn bị trước, tay chân thừa thãi không biết để đâu.

Rushia lay lay tay tôi, kéo tôi ra khỏi trạng thái bối rối bằng một cái chạm nhẹ nhàng. Em ấy cười với tôi, giọng nói thân thiện như thể chúng tôi đã là bạn lâu năm: ``Thôi nào, Kuon-senpai ơi!''

Marine giục tôi tham gia vào nhóm, không để tôi đứng một mình lạc lõng giữa bục giảng. Em ấy muốn tôi thực sự cảm nhận được sự gắn kết của nhóm Fantasy: ``Anh cũng phải làm theo bọn em chứ!''

``À\dots\ Hả?'' tôi ngượng ngùng vung tay lên cùng họ, cuối cùng cũng chấp nhận rằng mình đã bị cuốn vào câu chuyện này. Cả lớp bật cười trước vẻ ngoài ngơ ngác của tôi, không ai ngờ được rằng tôi cũng có ngày bị lôi vào một tình huống thế này.

Thế hệ thứ Ba hô vang lần cuối, gửi gắm tất cả tình cảm của mình vào câu nói cuối cùng. Mọi trò đùa nghịch đã biến mất, thay vào đó là sự chân thành từ đáy lòng. ``\textbf{Mãi mãi là những người bạn vĩnh cửu! Rushia-chan, chúng tôi mến bà/cậu lắm!}''

Cả lớp nhìn màn trình diễn đó mà ngơ ngác toàn tập, không ai có thể tin được vào những gì mình vừa chứng kiến trong buổi học cuối cùng. Đến cả tôi còn chẳng biết phải phản ứng ra sao trước sự ngượng ngùng này, đám bạn nhóm ``Rách'' cũng dở khóc dở cười như đang xem một bộ phim không đầu không cuối.

Tiếng vỗ tay cuối cùng cũng dần tan biến sau màn trình diễn đầy cảm xúc của năm cô gái, nhưng sự ấm áp mà họ để lại vẫn đang dậy sóng trong cả lớp học. Chúng tôi ngồi nhìn nhóm Thế hệ thứ Ba quây quần bên nhau, đủ mọi cảm xúc đan xen trong lòng. Có sự bối rối khó hiểu, có sự đồng cảm thấu hiểu, và hơn hết là\dots\ niềm hạnh phúc vô hình đang lan tỏa ra từng góc phòng.

Không ai trong chúng tôi kìm được câu thốt lên đầy ngạc nhiên, Kihataru là người đầu tiên phá vỡ sự im lặng: ``Cái quái gì thế này\dots''

Hakimi lắc đầu không dứt, mắt vẫn dán vào nhóm cô gái đang cười đùa trên bục: ``Mấy người này chuyển cảnh nhanh thật đấy\dots''

``Chịu luôn,'' tôi lắc đầu mỉm cười bất lực, không biết phải nói gì hơn.

Kijiru gật gù công nhận một điểm mà ai cũng thấy rõ, giọng nói trầm xuống chút: ``Nhưng phải thừa nhận là họ khăng khít thật sự nhỉ?''

``Ừ.'' Tôi đồng tình, gật đầu liên tục. ``Cái đó tôi cũng phải công nhận.''

Zetaro chêm vào một cách tự hào, như đang so sánh với chính nhóm của mình: ``Chắc cũng tầm cỡ anh em mình nhỉ?''

Hakimi lập tức bĩu môi phản đối, không đồng tình với sự so sánh đó: ``Tao không nghĩ chúng ta tình cảm sến sẩm được đến mức đấy đâu.''

Boji thở dài một tiếng, ánh mắt dõi theo nhóm cô gái với chút ngưỡng mộ: ``Công nhận tao cũng thấy hơi ghen tị với họ. Dám nói ra những điều thầm kín nhất với nhau trước mặt mọi người.''

Khan cũng gật gù đồng tình, ước mơ về một sự gắn kết tương tự cho nhóm của mình: ``Ước gì chúng ta cũng có thể gắn kết được như thế.''

Câu nói ấy như một lời nhắc nhở khiến tôi phải đối mặt với sự thật mà tôi đã cố gắng né tránh từ lâu. Thật lòng tôi không muốn nhắc đến chuyện đó trong buổi gặp mặt cuối cùng này. Nhưng sự thật thì mãi mãi vẫn là sự thật, không thể thay đổi.

``Tao nghĩ chúng ta có thể giữ được tình bạn lâu năm, thậm chí là rất thân,'' tôi nhìn năm cô gái với chút ngập ngừng, giọng nói trầm đi thấy rõ. ``Nhưng với họ\dots\ mong muốn đó cũng chỉ là một ước mơ hão huyền mà thôi.''

Hakimi lập tức ngơ ngác không hiểu, quay sang hỏi tôi với ánh mắt tò mò: ``Tại sao?''

Tôi thở dài một tiếng nặng nề, vị đắng của sự thật lan tỏa trong cổ họng: ``\dots Đây sẽ là lần cuối cùng họ có thể sum vầy đông đủ như thế này.''

Zetaro ngạc nhiên hẳn lên, nghiêng người về phía tôi như để nghe rõ hơn: ``Ủa? Sao lại thế?''

``Một người trong nhóm đó\dots\ trong tương lai gần thôi, sẽ chỉ còn là ký ức đối với những người còn lại. Em ấy sẽ phải rời khỏi nhóm, không thể tiếp tục cùng nhau nữa.''

Boji nhíu mày lo lắng, đầu tiên nghĩ đến những xung đột nội bộ thường thấy: ``Bộ nội bộ họ có xích mích gì à?''

``Không.'' Tôi lắc đầu dứt khoát, loại bỏ suy nghĩ đó. ``Là một tai nạn thì đúng hơn.''

Koheki đoán thử một cách thông thường, nghĩ đến những rủi ro đời thường: ``Kiểu tai nạn giao thông hay gì sao?''

``Cũng không phải. Là tai nạn nghề nghiệp,'' tôi sửa lại, giọng nói càng lúc càng trầm.

``Hả?'' Cả nhóm cùng thốt lên câu hỏi ngắn ngủi, không hiểu ý tôi là gì.

``Một tai nạn đủ nghiêm trọng để hủy hoại gần như toàn bộ sự nghiệp mà em ấy đã gầy dựng bao năm nay,'' tôi giải thích ngắn gọn, không muốn nói quá chi tiết.

Cả nhóm tôi cùng trố mắt kinh ngạc, không ngờ sự thật lại khắc nghiệt đến vậy: ``Phũ phàng đến thế cơ á!?''

Tôi lặng lẽ gật đầu, sự thật đó đã nằm trong lòng tôi từ lâu: ``Mọi người thừa biết nền công nghiệp idol ở bên đó khốc liệt đến mức nào rồi mà.''

Boji lặng người đi, dường như đã mường tượng ra được sự việc: ``Ồ\dots\ Tôi mờ mờ hiểu ra rồi.''

Zetaro ngập ngừng hỏi, vẫn còn chút mơ hồ về những gì mình từng chứng kiến: ``Vậy những gì tao thấy trên stream của họ\dots''

``\dots Đều có những đắng cay và áp lực phía sau,'' tôi cắt ngang câu hỏi của cậu bạn. ``Tao làm trong ngành tao biết rõ nhất. Nhiều khi chỉ một sai lầm nhỏ thôi là cút luôn. Nói thẳng ra là như vậy.''

Khan nhìn về phía bục giảng, ánh mắt dừng lại trên từng cô gái để đoán thử: ``Tôi đoán người đó\dots\ cũng vừa mắc phải sai lầm kiểu như thế đúng không?''

``\dots Chính xác.'' Tôi đáp giọng nhẹ hẫng, sự thật đó như một tảng đá đè nặng lên ngực.

Kihataru hạ thấp giọng hỏi, sợ rằng sẽ làm phiền đến cô gái kia nếu nghe thấy: ``Tao muốn hỏi\dots\ người đó là ai vậy?''

Tôi chỉ im lặng, không thể thốt ra cái tên đó ra giữa lớp học. Không cần nói ra, chỉ cần một cái gật đầu là đủ.

Hakimi nhìn tôi một lúc, rồi chậm rãi quay sang quan sát bục giảng, ánh mắt dừng lại trên cô gái đang cười tươi nhất: ``Tôi đoán là\dots\ ngôi sao chính của kịch bản vừa rồi\dots\ đúng chứ?''

Tôi im lặng gật đầu, xác nhận phỏng đoán của cậu bạn. Tất cả đám bạn tôi cùng quay sang nhìn Rushia, ánh mắt tràn ngập sự thương cảm không nói nên lời: ``Tội nghiệp bạn ấy ghê\dots''

Trái ngược hoàn toàn với bầu không khí ngậm ngùi của chúng tôi, năm cô gái trên bục giảng lúc này đang hồn nhiên cùng nhau cụng những ly trà sữa. Tiếng cười của họ vang lên trong veo, không biết về những gì đang chờ đợi ở phía trước.

Kijiru vội vàng chuyển chủ đề, cố gắng xua tan bầu không khí ảm đạm đang bao trùm nhóm mình: ``T-Thôi, dù sao hôm nay cũng là buổi học cuối cùng rồi, chúng ta cứ tận hưởng trọn vẹn khoảnh khắc này đi chứ?''

Mọi người nhìn nhau, đều nhận ra cậu ta nói rất có lý. Không nên để nỗi buồn làm hỏng buổi gặp mặt cuối cùng. ``Ừ.'' Tôi gật đầu đồng tình, cố gắng gạt bỏ những suy nghĩ tiêu cực ra khỏi đầu. ``Buồn ủ rũ lúc này cũng chẳng giải quyết được gì.''

Hakimi tặc lưỡi chấp nhận sự thật, không muốn suy nghĩ thêm làm gì: ``Bây giờ cứ xác định chuyện gì đến rồi cũng sẽ đến thôi.''

Zetaro lập tức hào hứng giơ ly trà sữa của mình lên cao, muốn thay đổi không khí: ``Anh em mình cũng zô một phát nhỉ!?''

Koheki lập tức phàn nàn, không hiểu cậu bạn hào hứng cái gì: ``Có phải uống bia méo đâu mà zô\dots''

Boji cười lớn vỗ vai cậu bạn Hóa học, không để ý đến những chi tiết nhỏ: ``Quan trọng gì, vui là chính!''

``Đó! \textbf{Anh em, nâng cốc lên!}'' Zetaro hô lớn, không để ý đến những lời phàn nàn.

``Tao đã bảo là có phải uống bia quái đâu cái thằng này!'' Koheki vừa chửi vừa vẫn giơ cốc lên, không muốn làm mất không khí.

Dù vẩn vơ vậy thôi, cả lũ chúng tôi vẫn ầm ĩ đập những ly trà sữa vào nhau, tiếng cụng cộp cộp vang vọng khắp góc lớp. Trông đám bạn tôi chẳng khác gì mấy ông bợm nhậu tụ tập ở quán bia hơi cuối phố, vừa ồn ào vừa lộn xộn. Nhưng khi ngước nhìn hội Thế hệ thứ Ba cũng đang hào hứng chẳng kém, tự dưng tôi nghĩ: hà cớ gì chúng ta phải thấy ngại ngần hay xấu hổ cơ chứ?

``Kuon-senpai! Lại đây với bọn em nào!'' Marine vẫy tay gọi lớn, đôi mắt sáng rực như vừa tìm được kho báu.

Rushia cũng hùa theo, nụ cười tươi rói nở trên môi khi hướng mắt về phía tôi. ``Nói chuyện với bọn em một chút đi chứ!''

``Ờ, rồi rồi, anh ra liền đây!'' tôi đáp lại nhanh gọn, vừa cười vừa xỏ chân vào đôi dép đi trong lớp.

Zetaro đẩy mạnh vào lưng tôi một cái, làm tôi suýt nữa thì văng ra trước. ``Ra nói chuyện với dàn em gái của mày đi kìa bạn, đừng có mà rụt rè.''

``Méo cần mày nhắc tao. Anh em cứ tự nhiên nhé, đừng có làm gì ngu là được!''

Từng bước đi ngang qua những dãy bàn còn vương lộn xộn mấy gói snack, tôi chợt nhận ra nhóm Thế hệ Thứ Ba - hololive Fantasy có một điều gì đó mà bản thân tôi mãi mãi cũng khó có được. Đó là sự gắn kết mật thiết - một điều tưởng chừng đơn giản nhưng lại là thứ quý giá nhất trong mọi mối quan hệ giữa người với người. Và có lẽ, trong số tất cả, Rushia-chan là người thấm thía điều đó hơn ai hết.

Cũng giống như tôi, em ấy thuộc tuýp người dễ dàng cảm thấy cô đơn giữa đám đông. Nhưng khác với tôi, em ấy lại không bao giờ chịu đầu hàng trước nỗi cô độc đó. Em ấy luôn tìm đủ mọi cách để mở lòng, để kết nối và kết bạn với thật nhiều người xung quanh.

Thật may mắn thay, số phận đã không đem lại cho em ấy những nỗi cô đơn đơn độc. Vào đúng lúc cần nhất, những người bạn đồng hành đã bước vào cuộc đời Rushia, lấp đầy những khoảng trống trong lòng em. ``Cô ấy may mắn thật sự,'' tôi thầm nghĩ trong đầu khi nhìn nhóm năm người đang cười nói ríu rít.

Họ cùng nhau tạo ra đủ thứ trò quậy phá khắp nơi, từ những buổi stream điên rồ đến những chuyến đi chơi xa đầy rẫy chuyện tiếu lâm. Mỗi khi có ai buồn hay gặp khó khăn, những người còn lại luôn là chỗ dựa đầu tiên để sẻ chia. Không có gì là quá lớn để họ không thể cùng nhau vượt qua.

Mối liên kết giữa họ bền chặt đến mức tưởng như không gì có thể chia cắt được, những sợi dây tình cảm vô hình đã thắt chặt năm trái tim lại làm một. Dù mai này có mỗi người một ngả, tình bạn ấy chắc chắn vẫn sẽ mãi tồn tại trong ký ức của mỗi người. Tôi thầm ao ước bản thân cũng có thể có được một sự gắn kết đẹp đẽ như vậy.

Nhưng\dots\ tôi cũng phải gửi tới em một lời xin lỗi.

Xin lỗi em nhé, Rushia.

Tôi, chú YAGOO, cùng toàn thể mọi người, dù đã cố gắng hết sức mình, nhưng\dots\ chúng tôi vẫn không thể cứu được em.

Sau này, khi chỉ còn lại một mình em đương đầu với những dông bão của cái gọi là ``dòng đời'', khi chẳng còn ai bên cạnh đồng hành nữa, con đường phía trước của em chắc chắn sẽ vô cùng chông gai và nghiệt ngã\dots

Nhưng dù vậy, anh vẫn hy vọng em giữ vững ngọn lửa nhiệt huyết ấy để vượt qua mọi thử thách. Và biết đâu đấy, nếu may mắn mỉm cười, một ngày nào đó chúng ta lại có thể gặp lại nhau.

Cảm ơn em, và tạm biệt nhé\dots

Uruha Rushia-chan.

%==========================================TRUYỆN PHỤ=========================================
\omk

\begin{center}
  5:00 sáng\\
  Nhà Hikawa - quận Kyuukami.
\end{center}

Tiếng chuông báo thức bật lên đúng năm giờ sáng, xé tan bầu không khí im lặng như tờ của căn phòng. Tôi chớp mắt mơ màng, quơ tay vội vàng chộp lấy chiếc điện thoại đặt trên bàn đầu giường rồi lóng ngóng quẹt tắt chuông.

Chậm chạp chống tay ngồi dậy, đầu óc tôi vẫn còn choáng váng như vừa lơ lửng trên mây. Tôi đảo mắt quanh căn phòng ngủ chật hẹp, trên chiếc giường rộng, em gái Naomi đang cuộn tròn ôm hai chú chó nhà là Mi và Tô, cả ba đang ngủ say sưa như chẳng có gì có thể đánh thức được. Còn ở dưới sàn nhà, bố tôi đã trải sẵn một chiếc đệm gấp để ngủ lại, trùm kín cổ bằng chiếc chăn bông và chỉ phát ra tiếng thở đều đặn qua mũi.

``Trời đã sáng rồi sao\dots?''

Tôi lẩm bẩm trong miệng, vừa vò rối mái tóc vừa gượng cơn buồn ngủ. Đúng là một giấc mơ kỳ lạ không thể nào tả nổi. Nhưng chẳng hiểu sao, toàn bộ những gì xảy ra trong mơ lại mang đến một cảm giác quen thuộc đến lạ kỳ, như thể tôi đã từng trải qua nó rồi.

Nói cho đúng thì sao nhỉ? Bỗng dưng trong mơ, tôi thấy bản thân quay trở lại mái trường cấp ba thân thuộc, đúng vào ngày học cuối cùng trước kỳ tốt nghiệp. Ở đó, tôi gặp lại nguyên đám bạn chí cốt của mình, và còn kỳ lạ hơn cả, cả năm cô gái của hội Fantasy cũng học chung lớp 12A3 với chúng tôi.

Thật là một giấc mơ quái thai đến không tưởng\dots\ Không lẽ đây là điềm báo cho một sự kiện sắp xảy ra? Tôi ngẩn ngơ nhìn vào màn hình điện thoại đang sáng dần, rồi lại thầm gật đầu với bản thân. Có lẽ đúng là vậy thật.

Tôi định kéo chân ra khỏi chăn để đi rửa mặt cho tỉnh táo, nhưng theo thói quen lâu năm, tôi vẫn chạm vào màn hình để kiểm tra tin nhắn trước khi làm gì cả.

Thứ Tư, ngày 23 tháng 2 năm 2022.

Một ngày làm việc bình thường như bao ngày khác trong lịch trình của tôi. Nhưng thứ khiến tâm trí tôi lập tức tỉnh táo hẳn không phải là dòng ngày tháng quen thuộc, mà là hai thông báo tin nhắn vừa nhảy lên màn hình, gửi đến chỉ vài phút trước.

Tin nhắn đầu tiên đến từ nhóm chat của hội ``Rách'' - nhóm bạn thân chí cốt từ hồi cấp ba của tôi. Mở cuộc trò chuyện ra, tôi không khỏi giật mình. Không chỉ vì tất cả bọn họ đều đã thức dậy từ rất sớm, mà là vì một điểm chung kỳ quặc đến rợn người\dots

Tất cả mọi người đều vừa trải qua đúng một giấc mơ giống hệt như tôi.

**Zetaro:** ``Ê, đêm qua t mơ thấy cái gì lạ lắm. Có ai bị như t ko (không)''

**Boji:** ``T cũng bị này''

**Hakimi:** ``What the hell man? Cả t cũng thế luôn''

**Koheki:** ``T cũng thế. Đm cái loz gì vậy''

Kihataru, Kijiru và Khan cũng nhanh chóng nhảy vào xác nhận. Tôi ngập ngừng gõ dòng tin nhắn trả lời:

**Kuon:** ``T cũng mơ như cm (chúng mày) này''

**Zetaro:** ``Ơ thế là ae (anh em) mình mơ cùng giấc mơ à''

**Khan:** ``Dễ là thế lắm chứ. T mơ thấy mình quay lại thời cấp 3\dots''

**Kihataru:** ``Xong rồi t còn có thể nói chuyện với ae nữa. Cảm xúc ấy lẫn lộn lắm\dots''

**Kijiru:** ``Đcm ảo thật''

**Zetaro:** ``Nhưng mà quan trọng nhất ấy\dots\ t còn mơ thấy năm đứa con gái chuyển vào lớp mình nữa cơ''

**Hakimi:** ``À à, t biết. Flare, Marine, Noel, Pekora\dots\ với cả Rushia? I guess?''

**Zetaro:** ``Chính chúng nó luôn! Cái này chắc thằng @Kuon biết rõ nhất này''

**Kuon:** ``À phải. T có nói trong giấc mơ là tao làm qlý (quản lý) mà''

**Kihataru:** ``Thật á''

**Kuon:** ``Cứ hỏi Zetaro-senpai là biết liền''

**Zetaro:** ``Nó nói đúng mà''

**Koheki:** ``vcl''

**Khan:** ``Thật luôn à Kuon?''

**Kuon:** ``T lừa ae bh (bao giờ) đâu!''

\dots

Và cứ như thế, cả lũ chúng tôi thao thao bất tuyệt xôn xao về giấc mơ nửa thật nửa ảo ấy.

Chúng tôi cùng nhau ôn lại từng chi tiết của buổi ``lễ tốt nghiệp'' tự phát trong giấc mơ hôm qua, cái buổi chia tay nhỏ nhắn mà chúng tôi chẳng ai dám chuẩn bị trước dành cho Rushia. Hệt như những gì vừa diễn ra trong mơ, cả nhóm đều dâng trào cùng một cảm xúc nghẹn ngào, như có một mảnh giấy nhúng chì đè nặng ở ngực. Nỗi thương xót ấy dành cho một người bạn sắp phải rời xa, một khoảng trống toang hoác trong tim mà chẳng ai dám nói ra, cũng chẳng bao giờ có thể lấp đầy được nữa.

``Thật sự éo dám tin là ae ta đều mơ thấy cùng một cảnh tượng'' Zetaro gõ dòng tin nhắn chậm rãi, như đang sợ hãi với sự trùng hợp quái lạ ấy. Không ai trong nhóm nhắn thêm gì nữa, chỉ có những dấu chấm ``đang soạn tin'' nhấp nháy liên tục, thay cho những câu nói chẳng biết phải bắt đầu từ đâu.

Mọi người đều cùng nghĩ về một điều: nhóm anh em ``Rách'' của chúng tôi vẫn còn quá may mắn khi có thể ở bên nhau mãi, dù đi đâu cũng có thể gặp lại. Nhưng với nhóm năm cô gái ấy\dots\ nếu vắng bóng Rushia-chan, mảnh ghép Fantasy sẽ vĩnh viễn chẳng bao giờ trọn vẹn được nữa.

Cuộc trò chuyện im lặng kéo dài mãi, những dòng tin nhắn thưa dần cho đến khi đồng hồ điện thoại điểm đúng 6 giờ sáng. Tôi tắt điện thoại đi, đặt lên bàn đầu giường, rồi ngước mắt nhìn ra ngoài cửa sổ phòng mình. Bầu trời khu vực Kawachi hôm nay xám xịt như được phủ một lớp tro nguội, u tối đến nghẹt thở, đúng như cái không khí thường có của những ngày mưa dông cuối đông, khi cái lạnh vẫn còn bám rít lấy từng góc phố.

Bất chợt, một vệt màu xanh lẻ loi bay ngang qua khung cửa sổ, kéo tầm mắt tôi theo sau. Đó là một chú bướm nhỏ, nó chầm chậm nghiêng cánh rồi đậu nhẹ nhàng xuống mặt kính lạnh. Nó đứng lặng yên ở đó rất lâu, cánh khép lại như đang giữ một bí mật, tưởng chừng như đang muốn nhắn gửi điều gì đó đến tôi trong căn phòng này.

Con bướm ấy\dots\ tại sao lại nhìn quen thuộc đến thế nhỉ?

Tôi khẽ đứng dậy, không dám di chuyển mạnh sợ làm nó bay mất, mắt dán chặt vào đôi cánh nhỏ bé. Nó khoác lên mình sắc xanh cánh chả rực rỡ, xen lẫn những đường gân đen vắt ngang, và vài đốm trắng li ti ở rìa viền cánh. Đấy chính là loài bướm Pareronia boebera, một giống bướm mà tôi chỉ từng thấy trong sách, đặc trưng của vùng Philippines xa xôi.

Không lẽ nào\dots\ nó đến từ nơi đó?

Tôi đứng bất động ở cửa sổ, đầu óc trôi đi miên man giữa bao nhiêu câu hỏi chưa có lời đáp. Trong cái khoảnh khắc tôi chới với giữa thực tế và giấc mơ vừa qua, chú bướm đã vỗ cánh nhẹ nhàng bay mất từ lúc nào. Nó biến mất vào màn mưa sắp rơi, để lại trong tâm trí tôi vô vàn thắc mắc ngổn ngang, chẳng thể nào gói gọn lại được.

Chẳng nhẽ\dots\ chuyện mà tôi đã sợ hãi hàng tháng trời, đã thực sự xảy đến rồi sao? Mặc dù bản thân tôi biết rõ mình đã chuẩn bị tinh thần cho tình huống xấu nhất này\dots

\dots

Thở dài một tiếng, tôi chợt nhớ ra: hình như mình vẫn còn một tin nhắn chưa đọc thì phải?

Tin nhắn còn lại đến từ A-san. Hẳn là chị ấy phải đang gặp chuyện gì đó vô cùng gấp gáp nên mới gửi cho tôi vào lúc nửa đêm như thế, chính xác là từ tận mười một giờ đêm qua. Bình thường tôi hay thức muộn, nhưng tối qua vì xong việc sớm nên tôi đã đi ngủ từ lúc mười giờ.

Thế nhưng, những dòng chữ hiện ra trên màn hình điện thoại lại khiến tim tôi chết đứng.

Điều tồi tệ nhất\dots\ đã thực sự xảy đến rồi.

``Xin lỗi nhé, Kuon-san. Nhưng chị có thể nhờ nhóm Houshou-san một việc được không? Uruha-san đã mất tích kể từ khi em ấy bị sa thải từ tối qua rồi\dots''

\begin{center}
  {\abv Giải thích}
\end{center}

Phần đầu tiên dựa trên hoàn toàn theo trí nhớ (cộng với sự tạp nham của các tình huống khác) của người dịch khi còn là học sinh cấp ba.

Nửa đầu phần thứ hai thì dựa theo cliché anime được lụm nhặt trên mạng.

Nửa sau phần thứ hai, như đã nói ở trên, dựa theo trích đoạn trên stream của Uruha Rushia, tên 【１００万人thanks】振り返り(地獄)をしながら凸待ちをする(重大発表有)【潤羽るしあ/ホロライブ】 (hiện đã bị ẩn do kênh không còn hoạt động).

Về nội dung đoạn trích, đó đơn giản chỉ là cảm nghĩ của các thành viên Gen 3 về Rushia khi cô ấy đạt 1 triệu lượt đăng ký (đó cũng là mục đích để cô ấy làm buổi stream ấy). Người dịch đã mượn ý tưởng này để biến tấu đi sao cho giống nhất với một buổi tốt nghiệp (mà chính Rushia cũng nói nó khá giống như vậy). Đoạn trích lấy từ video của @momonofuen1427.

Tên ``Anwa'' lấy từ chữ Hán 安和, dịch ra là ``Yên Hòa''. Đó cũng là tên trường THPT (cấp ba) mà người dịch đã từng học.

\end{document}